% Mechanistic Validity: A Validity Theory, Research Methodology, and Evidence Standard for Mechanistic Claims
\pdfoutput=1

\documentclass{article}

\usepackage[preprint]{tmlr}

\newcommand{\openreview}{\url{https://openreview.net/forum?id=XXXX}}
\def\month{06}
\def\year{2026}

\usepackage{amsmath,amssymb}
\usepackage{booktabs}
\usepackage{graphicx}
\graphicspath{{figures/}}
\usepackage{xcolor}
\usepackage{hyperref}
\usepackage{cleveref}
\usepackage{enumitem}
\usepackage{placeins}
\usepackage{float}
\usepackage{multirow}
\usepackage{array}
\usepackage{longtable}

% --- Macros ---
% No abbreviation — spell out "Mechanistic Validity" in full
\newcommand{\ie}{i.e.}
\newcommand{\eg}{e.g.}
\newcommand{\cf}{cf.}

% tight tables
\newcommand{\tighttable}{\setlength{\tabcolsep}{4pt}\renewcommand{\arraystretch}{1.05}}

% add space between caption and table rule
\usepackage{caption}
\captionsetup[table]{skip=8pt}

\title{Mechanistic Validity: A Validity Theory, Research Methodology, and Evidence Standard for Mechanistic Claims}

\author{\name Elliot Tower \\
      \email elliot@elliottower.ai}




\begin{document}

\maketitle

% ============================================================
\begin{abstract}
Mechanistic claims are ubiquitous---\emph{this gene causes this disease}, \emph{this pathway mediates this outcome}, \emph{this circuit implements this computation}. What makes such a claim valid?

We argue that a mechanistic claim is only as strong as its weakest validity dimension. Five validity types---construct, measurement, internal, external, and interpretive---form a dependency chain: construct validity is a ceiling on all others, so no amount of causal evidence rescues a claim whose target concept is incoherent. This is a theory of validity, not a taxonomy: it predicts that claims will cap at the tier corresponding to their weakest dimension, and that single high-scoring metrics cannot compensate for untested dimensions.

We operationalize this theory as Mechanistic Validity, drawing on philosophy of science, measurement theory, and validation methodology from neuroscience, pharmacology, and genetics. The framework evaluates; the methodology produces what is evaluated. For practitioners, the realistic ask is three questions: Is the description mode declared? Is the claimed tier supported by evidence? Does the evidence tier match the claim tier?

We audit sixteen mechanistic claims drawn from fifteen papers: sparse-autoencoder features, probing classifiers, steering vectors, and the global workspace hypothesis. The framework discriminates, and it separates claims that share a paper. Induction heads reach Triangulated for in-context token copying and Disconfirmed for general in-context learning, on evidence the original authors named as their own refutation condition. The IOI circuit reaches Causally Suggestive under its most lenient reading: rival head sets leave that reading standing, and what caps the claim is method disagreement, its headline faithfulness number spanning below 0\% to over 100\% depending on the ablation protocol. We propose evidence cards as a standard reporting unit for mechanistic claims, making gaps visible and the path to closing them explicit.
\end{abstract}

\newpage
{\small\setlength{\parskip}{0pt}\setcounter{tocdepth}{3}\tableofcontents}
\newpage

% ============================================================
\section{Introduction}
\label{sec:intro}

Mechanistic interpretability has produced a remarkable catalog of internal mechanisms:
induction heads that implement in-context copying \citep{olsson2022context}, a multi-role
circuit for indirect object identification \citep{wang2023interpretability}, a greater-than
algorithm \citep{hanna2023greater}, and sparse autoencoder features proposed to correspond
to human-interpretable concepts \citep{cunningham2023sparse, bricken2024monosemanticity}. The field has also begun to
build evaluation infrastructure: faithfulness metrics for circuits
\citep{miller2024transformer}, the Mechanistic Interpretability Benchmark for comparing
localization methods across tasks and models \citep{mueller2025mib}, SAEBench for auditing
sparse autoencoder quality \citep{karvonen2025saebench}, and causal scrubbing for testing
mechanistic hypotheses via behavior-preserving resampling \citep{chan2022causal}.

But how well-supported are these claims? Across the field's most prominent results, a pattern recurs: strong evidence along one dimension, untested assumptions along others. Faithfulness scores that are method-conditional. Metrics that saturate on random models. Constructs that cannot be separated from neighboring computations. Specificity controls that were never run. Labels that imply more than the evidence supports. Each is a gap in the \emph{kind} of evidence collected, which no amount of execution quality would close, and \Cref{sec:mi-failures} catalogs them systematically. The pattern is not unique to MI: every field that makes mechanistic claims has encountered the same structural problems (\Cref{sec:motivation}). Shared metrics are not yet shared validity criteria.

Mechanistic Validity is a structured framework that takes the validity vocabulary the field already uses informally and makes it operational, falsifiable, and auditable. It evaluates claims through a seven-layer pipeline: a \emph{description mode} declares what kind of explanation is being offered. \emph{Evidence families} classify the sources of signal. \emph{Metrics} are concrete tests that measure the claim. \emph{Criteria} test whether the evidence meets the required conditions. Five \emph{validity types}---construct, measurement, internal, external, and interpretive---determine which dimensions the evidence addresses, in dependency order. \emph{Synthesis protocols} aggregate evidence and assign \emph{verdicts}: structured validity profiles and tier assignments, from Proposed through Causally Suggestive, Mechanistically Supported, and Triangulated, to Validated.

The pipeline has two uses. When \emph{running experiments}, evidence builds upward from metrics through criteria to a verdict. When \emph{auditing a claim}, the process reverses: starting from a verdict, the framework identifies which criteria are unmet and whether the description mode is supported by the available evidence.

The same criteria apply to circuits, SAE features, probing classifiers, steering vectors, transcoders, and crosscoders---the questions ``is the construct well-defined?'' and ``is the evidence causal?'' do not depend on how the evidence was produced.

Its unit of analysis is a single causal mechanistic claim---not a model, a method, or a benchmark score---and the evidence a claim owes is fixed by the description mode it declares. We describe lessons from validity failures across sciences (\Cref{sec:motivation}), the framework (\Cref{sec:framework}), the concrete methodology for generating evidence (\Cref{sec:methodology}), its theoretical foundations (\Cref{sec:foundations}), its application to sixteen published claims spanning circuits, features, steering vectors, and a subspace-level result (\Cref{sec:case-studies}), and its implications for the field (\Cref{sec:discussion}).


% ============================================================
\section{Motivation: What Happens Without Validity Criteria}
\label{sec:motivation}

Every field that makes mechanistic claims has encountered the same structural problem: strong individual results that do not add up to valid conclusions. The specific failure modes differ, but the pattern is the same---an untested link in the inference chain from measurement to claim. We catalog these failures across sciences, machine learning, and mechanistic interpretability, organized by the validity criterion each violates.

\subsection{Validity failures across the sciences}

\paragraph{Dead salmon fMRI (M2: baseline separation).}
\citet{bennett2010neural} scanned a dead Atlantic salmon while presenting it with photographs of humans in social situations. Using standard uncorrected thresholds, 16 significant voxels appeared in the salmon's brain cavity; with appropriate multiple comparisons corrections, none survived. \citet{eklund2016cluster} extended this systematically, finding that standard fMRI software produced false positive rates up to 70\%. This same structure can be seen in MI, if a method measures faithfulness without testing against a size-matched random subnetwork.

\paragraph{The unreported denominator in MI (M7: selection correction).}


\paragraph{Cardiac stents for stable angina (I4: specificity).}
For three decades, stenting was standard treatment for stable angina (~500,000 procedures/year in the US). The ORBITA trial \citep{alhabib2018orbita} was the first placebo-controlled test: half of patients received a sham procedure with no stent placed. The result: no statistically significant difference in exercise tolerance or angina symptoms. The entire symptomatic benefit was placebo. The effect was real. The causal attribution was not---it had never been tested against a placebo control. Circuit ablation without a size-matched random baseline faces the same problem.

\paragraph{Candidate gene psychiatry (I7: confound control).}
The candidate gene era (1990s--2010s) produced thousands of associations between gene variants and depression risk. \citet{border2019no} tested the 18 most-studied candidate genes across 621,214 individuals: the candidate genes ``(with the possible exception of DRD2) were no more associated with depression phenotypes than genes chosen at random.'' Twenty-five years of literature produced real associations in the samples---but population stratification and selective reporting were never ruled out. When \citet{border2019no} controlled for these confounders at scale, the signal disappeared. This is the same evidence pattern that appears in MI probing studies: a positive result that survives only as long as the confounds go unmeasured.

\paragraph{Cold fusion (M1/M2: measurement reliability and baseline separation).}
In 1989, Pons and Fleischmann announced room-temperature nuclear fusion. Several labs initially confirmed the result; within months, all retracted. Neutron detectors gave false positives when exposed to heat, and the original paper reported a gamma peak without its Compton edge---a physical impossibility indicating instrumentation error. The instrument was producing the expected signal for the wrong reason. Validating the measurement independently of the hypothesis---testing whether the detector fires under conditions where fusion cannot occur---would have killed the claim in weeks, not years. In MI, this is equivalent to a metric never tested on a null model---it cannot distinguish the mechanism from its own response properties.

\paragraph{BOLD signal (M1: measurement reliability).}
Functional MRI does not measure neural activity---it measures blood oxygenation changes coupled to neural activity through a hemodynamic mechanism that varies by brain region, age, and individual. Every fMRI claim that says ``region X activates during task Y'' is technically a claim about hemodynamics, not neurons---and the two can dissociate. The instrument measures a proxy, not the target. Activation patching has the same structure: it measures the joint effect of the component's causal role and the network's compensatory response, and a single experiment cannot decompose the two.

\subsection{Validity failures in machine learning}

\paragraph{Emergent abilities mirage (M4: calibration).}
\citet{schaeffer2024emergent} showed that ``emergent abilities'' in large language models---apparent phase transitions where capabilities appear suddenly at scale---are artifacts of nonlinear evaluation metrics. Exact-match accuracy creates a sharp threshold where the model goes from 0\% to near-100\%; switch to linear metrics (token edit distance, Brier score) and the ``emergence'' disappears---the ability was always there, growing smoothly. M4 requires that reported numbers correspond to meaningful quantities. The metric created the phenomenon: emergence was real in the measurement but not in the model.

\paragraph{COVID-19 ML diagnostics (I7: confound control).}
\citet{roberts2021common} reviewed 61 studies claiming to detect COVID-19 from chest X-rays and CT scans. Not a single model was clinically usable. \citet{degrave2021ai} showed why: saliency maps revealed that models relied on laterality markers, image border artifacts, and patient positioning---not lung pathology. I7 requires that confounders be identified and controlled before a causal claim is made.

\paragraph{Google Flu Trends (C1/M3: construct definition and measurement stability).}
\citet{lazer2014parable} document a model that tracked the wrong construct and then drifted away from it. The fitting procedure selected the best matches among 50 million search terms against 1{,}152 data points, and the developers report weeding out terms that correlated strongly with the reference data while having nothing to do with influenza, ``such as those regarding high school basketball.'' The resulting instrument was, in their phrase, ``part flu detector, part winter detector''---which is why it missed the non-seasonal 2009 H1N1 pandemic entirely, and why it later overshot the reference measure in 100 of 108 weeks, at one point predicting more than double the observed rate. A second and distinct failure compounded it: Google changed its own search algorithm and interface repeatedly, so the signal shifted underneath a measure that had no stability check. The authors state the lesson in terms this framework adopts wholesale---``quantity of data does not mean that one can ignore foundational issues of measurement and construct validity and reliability.''

\paragraph{Shortcut learning (E2/C4: generalization and discriminant validity).}
\citet{geirhos2019imagenet} showed that ImageNet-trained CNNs are primarily texture detectors, not shape detectors---directly opposite to human perception and to what the field assumed. A cat-shaped object with elephant texture is classified as an elephant. Models achieving 90\%+ accuracy on the benchmark had learned a fundamentally different construct than intended. E2 requires that findings generalize across conditions; C4 requires that the measure distinguish the target construct from its neighbors.


\paragraph{Disagreement among attribution methods (C3: convergent validity).}
\citet{krishna2022disagreement} run six post-hoc explanation methods over six predictive models on four datasets and report that the methods ``often disagree in terms of the explanations they output.'' The second finding is the one that matters here: practitioners ``often employ ad hoc heuristics when resolving such disagreements,'' having no principled account of what agreement between two attribution methods would establish. Convergent validity (C3) states the condition the field is missing---evidence from sources that fail differently should agree, and where it does not, the disagreement is information about the instruments rather than a tie to be broken by preference. The case is instructive for what it exposes: six methods disagreeing is a result about the instruments, and the field had no framework that could read it as one.

\subsection{Validity failures in mechanistic interpretability}
\label{sec:mi-failures}

\paragraph{SAE metric baselines (M2: baseline separation).}
\citet{karvonen2025saebench} run a random-initialization control on five of their eight SAE evaluation metrics, and one inverts: Targeted Probe Perturbation scores \emph{higher} on a randomly initialized model than on a trained one. The authors read this as incomparability rather than failure --- the metric ``measures relative changes in probe accuracy within a given model, not across distinct models'' --- which is a construct problem rather than a measurement one, and is the harder of the two to notice. The metrics are reliable---they give the same answer when rerun---but a high score does not distinguish genuine feature quality from chance. Without a null distribution, there is no way to know whether an SAE feature score reflects real structure or noise that the metric cannot filter.

\paragraph{Vacuous nonlinear IIA (M2: baseline separation).}
\citet{sutter2025nonlinear} show that unconstrained nonlinear alignment maps achieve 100\% interchange-intervention accuracy on randomly initialized models. The IIA metric is working correctly; it is measuring the alignment map's flexibility, not the model's representational structure. Structured methods with identifiability constraints do not show this failure. The lesson is that a metric without a baseline is a metric without meaning: perfect scores are not evidence if random models also achieve them.

\paragraph{One control, three failures (M2: baseline separation).}
These are not three metrics with three problems. \citet{heap2025sparse} complete the set:
automated interpretability scores on sparse autoencoders trained over \emph{randomly
initialized} transformers are similar to those from trained models, so the instrument does
not distinguish model structure from dictionary structure. Across two method families ---
sparse autoencoder evaluation and causal abstraction --- a reported number turns out to be
one a random model earns as readily as a trained one, and in the SAEBench case more
readily. That a single control catches all three is what makes baseline separation a
criterion rather than a courtesy. It does not test any particular metric. It tests whether
a reported quantity is about the model at all, and a field can run every other check on
this list and still not have asked that question.

\paragraph{IOI circuit specificity (I4: specificity).}
The IOI circuit \citep{wang2023interpretability} demonstrates necessity and sufficiency---ablating the circuit degrades IOI performance, and isolating it recovers 87\% of the logit difference. Specificity is a separate question, and the evidence that exists is negative. \citet{wu2026demystifying} test transfer \emph{within} the task, across prompt templates: per-sample circuits of 712 and 754 edges share a Jaccard index of 0.126, and patching one sample with another's circuit moves the probability difference from $+0.120$ to $-0.224$---the transferred circuit drives behavior in the direction opposite to the intended intervention. Across the task as a whole, GPT-2 small on subject--verb agreement yields 826 edges per sample circuit against a union of 8{,}735. ``The circuit for task $T$'' is not a single object at the sample level. Cross-\emph{task} specificity has been tested and the answers disagree by model and by quantity---78\% head overlap with the Colored Objects circuit in GPT-2 Medium \citep{merullo2024circuit} against near-zero cross-task faithfulness in GPT-2 small \citep{hanna2024faith}---so the within-task failure is the cleaner result. Specificity requires establishing what the circuit does \emph{not} do, and here what it does varies with the prompt.

\paragraph{Docstring circuit negative controls (I4: specificity).}
The docstring circuit \citep{heimersheim2023docstring} is established by resample ablation on a 4-layer model that also handles natural language and general Python, and the effect of that circuit on any non-docstring behavior is never measured. Only predictions that should pass are tested. The gap is narrow and worth stating precisely: the authors control the rival algorithm carefully by prompt design, so what is missing is the converse test---whether removing the circuit costs the model anything it was not built to explain.

\paragraph{Knowledge neuron off-target effects (I4: specificity).}
\citet{dai2022knowledge} do measure off-target effects, and their own numbers point both ways. Erasing a relation raises that relation's perplexity by 36.7--141.2\% while raising other relations' by 1.1--10.1\%, which supports specificity. Their fact-updating experiment does not: editing the identified neurons raises inter-relation perplexity by 7.2 against 4.3 for \emph{random} neurons, so the located units damage unrelated facts more than arbitrary ones do. The paper prints the number and reads the table as showing ``little negative influence on other knowledge predictions.'' The off-target effect was measured; what failed was the reading of it.

\paragraph{Circuit non-uniqueness (I5: rival mechanism exclusion).}
\citet{meloux2025everything} establish non-identifiability by exhaustive enumeration on Boolean functions over small multi-layer perceptrons, where every candidate explanation can be checked: multiple circuits replicate the behavior, and multiple interpretations fit each circuit. The same holds where enumeration is impossible. \citet{wang2023interpretability} report that their own na\"ive circuit reaches comparable faithfulness, and \citet{chen2026circuits} exhibit two IOI circuits at perfect accuracy sharing 4.1\% of their edges, close to chance overlap. The ``IOI circuit'' is \emph{one of several} sufficient subnetworks, not \emph{the} mechanism. This is Quine's underdetermination problem \citep{quine1951two} made concrete: the evidence is consistent with multiple incompatible circuits, and the standard procedure cannot distinguish between them. Without a method for resolving between rival circuits, faithfulness scores select one of many equivalent explanations.

\paragraph{Method-conditional faithfulness (E1: intervention reach).}
\citet{miller2024transformer} vary six methodological choices in how the IOI circuit's faithfulness is computed, and the same quantity ranges from below 0\% to well over 100\%. The headline 87\% is one point in that range, and the ordering is not even stable: resample ablation scores systematically lower than mean ablation at the node level, while the edge-level, specific-token-position variant that ``best represents the hypothesis of Wang et al.'' has a median well over 100\%. The causal claim is method-conditional in both directions, which is worse than a claim that merely degrades under a stricter test. A finding whose sign depends on the intervention protocol has not established external validity.

\paragraph{Gender bias circuits (C4: discriminant validity).}
The construct ``gender bias circuit'' presupposes that bias separates from gender competence---that one could intervene on the biased prediction without disturbing correct pronoun resolution. \citet{vig2020gender} localize bias mediators by causal mediation analysis and never localize a competence mechanism, so no overlap statistic exists, and the definitionally gendered items that would carry the test are excluded from the measurement by design. Discriminant validity is untested rather than failed, which is the more common and the more easily missed condition: a claim whose central presupposition has never been measured looks, from the outside, exactly like one that has passed.

\paragraph{Othello ``world model'' (V4: interpretive inflation).}
\citet{li2023othello} recover Othello board state with nonlinear probes and show by 2,000 interventions that the representation is causally used for move prediction---the evidence is stronger than the label's critics usually allow. What is missing is a contrast: the paper defines ``world model'' as an understandable model of the process generating the sequences, and then runs no experiment separating that from a decodable state summary the network happens to consult. The gap runs between evidence for one reading and a name that asserts a second. \citet{vafa2024evaluating} later build the discriminating instrument, and the outcome splits the paper's own two checkpoints.

\subsection{Common structure}

\Cref{tab:cross-science} maps these failures to Mechanistic Validity criteria alongside the MI cases from \Cref{sec:mi-failures}.

\begin{table}[H]
\centering
\small
\caption{Validity failures across fields, mapped to Mechanistic Validity criteria. Each failure involves a strong individual result with an untested link in the inference chain.}
\label{tab:cross-science}
\renewcommand{\arraystretch}{1.15}
\begin{tabular}{@{}l>{\raggedright\arraybackslash}p{2.0cm}>{\raggedright\arraybackslash}p{2.5cm}>{\raggedright\arraybackslash}p{5.5cm}@{}}
\toprule
\textbf{Case} & \textbf{Field} & \textbf{Criterion} & \textbf{Untested link} \\
\midrule
Dead salmon fMRI & Neuroscience & M2 Baseline sep. & No null distribution; 70\% false positive rate \\
Cold fusion & Physics & M1/M2 Reliability & Instrument error; \$100M spent \\
BOLD signal & Neuroscience & M1 Reliability & Hemodynamics $\neq$ neural activity \\
SAEBench metrics & MI & M2 Baseline sep. & One metric scores \emph{higher} on a random-init model \\
Sutter IIA & MI & M2 Baseline sep. & 100\% IIA on random models \\
\midrule
Cardiac stents & Medicine & I4 Specificity & Sham-controlled: zero benefit \\
Knowledge neurons & MI & I4 Specificity & Fact updating hurts unrelated relations more than random neurons do \\
Candidate genes & Genetics & I7 Confound ctrl & 18 genes $=$ random in $n=620$K \\
COVID-19 ML & ML & I7 Confound ctrl & 0/61 models clinically usable (methodological flaws and bias) \\
\midrule
IOI circuit & MI & E1 Interv.\ reach & Same quantity spans below 0\% to over 100\% across methods \\
Shortcut learning & ML & E2/C4 Generalization & Texture $\neq$ shape understanding \\
\midrule
Emergent abilities & ML & M4 Calibration & Metric created the phase transition \\
Othello ``world model'' & MI & V4 Interp.\ inflation & Label never contrasted with ``causally-used state summary'' \\
\bottomrule
\end{tabular}
\end{table}

Despite the diversity of these failures---across physics, neuroscience, genetics, medicine, and machine learning---they share a common structure: in each case, an individual experimental result was strong, but the \emph{inference chain} from measurement to conclusion had an untested link. The result was real; the conclusion was not warranted by the evidence pattern. Claims migrated upward in inferential scope faster than supporting validity criteria could be established: a significant voxel cluster became a brain region that responds to social stimuli, a gene-environment interaction in 200 participants became a biomarker for depression, a high accuracy score became visual understanding, a sharp metric curve became an emergent ability, and symptom relief after stenting became evidence of a causal mechanism.

Each failure mode maps onto a specific criterion: measurement unreliability onto M1, lack of baselines onto M2, missing specificity controls onto I4, method-conditional results onto E1, interpretive inflation onto V4, unfalsifiable constructs onto C1. The framework makes the distance between evidence and claim visible, so that the field can prioritize the experiments that would close the gap.


% ============================================================
\section{Related Work}
\label{sec:related}

This section situates each of the paper's four claims against the literature that bears on it. The criteria are almost entirely imported and we name their sources; what is new is the object being assessed, the treatment of a claim as a family of readings, three interpretive criteria, and the separation of completeness from evidential support.

% ------------------------------------------------------------------
\subsection{The unit of analysis}
\label{sec:related-unit}

\paragraph{What this framework takes as given.} A mechanistic claim rests on a prior answer to what kind of thing a mechanism is, and that answer determines when two descriptions pick out the same mechanism and what evidence can bear on either. \citet{tower2026mechviews} sets out nine such views---object, role, subspace, structural, process, contrastive, perspectival, instrumental and stratified---each with an identity criterion, and shows that disagreements in the field frequently reduce to unstated mismatches between them. We take a claim's view as an input. Everything below applies within one, and the restriction is ours rather than \citeauthor{tower2026mechviews}'s: comparing hypotheses requires scoring each against the same evidence, and an evidence item that bears on a set of components need not bear at all on a claim about a subspace. \citeauthor{tower2026mechviews} does compare across views, and treats convergence between them as stronger evidence than any single view supplies, but that comparison runs between results already established within each view rather than between hypotheses competing for the same evidence.

That division is drawn in the same place from the other side. \citeauthor{tower2026mechviews} identifies a residue that view specification cannot dissolve: two researchers holding the same view, the same identity criterion and the same evidence type can still disagree about a claim. What remains after the ontology is fixed is which rival mechanism within that view the evidence favours, and how strong a reading of the winner it will carry. Those are the two questions this framework answers.

Validity frameworks in wide use fix the proposition under assessment before assessing it, because their fields supply one. \citet{cronbach1955construct} and \citet{messick1995validity} developed construct validity for measurement instruments: the object is a test score, and the question is whether it means what its user takes it to mean. \citeauthor{campbell1963experimental}'s taxonomy, extended by \citet{shadish2002experimental}, evaluates causal inference from experiments: the object is a study, and the question is whether its design licenses the conclusion drawn. \citet{guyatt2008grade} rate the certainty of a body of evidence bearing on one clinical question, reporting a level together with the domains that lowered it. The US Preventive Services Task Force recommends for or against a stated intervention in a stated population and issues an \emph{I statement} when the evidence cannot support either \citep{uspstf2021}. In each case the proposition is determined in advance, and validity is a judgment about how well evidence supports it.

None of these objects is a claim about how a system works. A test score, a study design, a treatment effect and a gene--disease assertion each have a settled propositional content that an assessor can take as given. A mechanistic claim has parts, an organization, a level of description, and a scope, and each of those can be asserted at more or less strength independently of the others. Those four dimensions make it a different object, not a special case of the others, and they are why the framework needs machinery the others do not: description modes, an interpretive validity block, and a procedure for identifying which proposition is being assessed before assessing it.

Assurance cases are the nearest infrastructure in another direction. Claims-Arguments-Evidence and Goal Structuring Notation decompose a top-level claim into subclaims supported by evidence, which is structurally similar to what an audit produces. They differ in purpose: an assurance case is built by a system's proponent to argue that a specific claim holds, and its confidence measures ask how much the argument's structure justifies belief. An audit is built by a third party against a claim already published, and must handle the case where the proponent's own statement of the claim is ambiguous.

% ------------------------------------------------------------------
\subsection{Readings, and the evidential frontier}
\label{sec:related-readings}

The observation that claims carry variable strength is old, and the vocabulary for it is well developed. \citet{toulmin1958uses} makes the \emph{qualifier} a component of argument structure precisely so that a claim's force is stated rather than assumed. The hedging literature in philosophy of science studies the same phenomenon from the author's side, analyzing when narrowing a claim's scope legitimately increases its warrant. \citet{niiniluoto1987truthlikeness}, developing Popper, replaces a binary confirmed/disconfirmed verdict with graded truthlikeness across theories. Each of these treats strength as a scalar attached to a proposition, or as a comparison between distinct theories. None decomposes one claim into an ordered family of propositions and scores each against the same evidence.

Two nearer neighbours come from evidence-based medicine, and one of them is close enough that we adopt it as precedent rather than distinguish it. The surrogate-outcome literature scores the relationship between a surrogate and a patient-centred outcome against three explicitly separate levels of evidence: Level~3, biological plausibility from pathophysiological reasoning; Level~2, a consistent association from observational or interventional studies; Level~1, evidence that treatment-induced changes in the surrogate correspond to treatment-induced changes in the final outcome \citep{euhta2024outcomes}. These are three different propositions about one relationship rather than three quality grades of one proposition, satisfying a weaker level establishes nothing about a stronger one, and papers are routinely criticised for conflating them. That is the structure this paper generalizes. It differs from ours in that its levels are fixed by domain---every surrogate claim is scored against the same ladder---whereas readings are derived per claim from what the paper asserted, and vary along two axes: mechanistic specificity and scope.

\citet{hultcrantz2017grade} report the same diagnosis from inside GRADE: a certainty rating is uninterpretable until the target of the rating is specified, because reviewers routinely rate certainty without stating which magnitude of effect they mean. GRADE's remedy is to require authors to declare one target per outcome. Ours is to score several and report which survive, because for a published mechanistic claim the target is often genuinely plural---the paper asserts one reading in its title and another in its body---and forcing a single choice hides the discrepancy that matters.

Two further literatures explain where the ambiguity comes from and how it spreads. The Duhem--Quine underdetermination tradition establishes that one evidence set is compatible with several theories given different background assumptions, which is the general condition of which our readings are a controlled and explicit case. And \citet{greenberg2009citation} supplies the empirical mechanism by which readings drift: tracing a citation network, he documents \emph{citation transmutation}, in which a hypothesis becomes an asserted fact through citation alone, with supportive citation paths multiplying without any new primary data. This is what produces the gap between what a paper claimed and what the field takes it to have shown, and it is why an audit that scores only the field's reading returns a verdict on authors who stated their limits carefully.

\citet{platt1964strong} decomposes a different axis, and the two axes are easily conflated. Strong inference proceeds by devising ``a crucial experiment (or several of them), with alternative possible outcomes, each of which will, as nearly as possible, exclude one or more'' hypotheses. Platt disaggregates the space of \emph{rival} mechanisms; readings disaggregate one mechanism claim into propositions of increasing strength. Both decompositions appear in this framework and they are orthogonal: the rival axis is what I5 tests and what the Underdetermined label records, and the reading axis is what the tier ladder measures. \citet{colombo2015models} formulate the same pair in defining a how-actually model as one that is both well confirmed by the available evidence and best confirmed among the available models of the phenomenon. The first condition is the tier; the second is rival exclusion. This also explains why Underdetermined terminates rather than caps: resolving it requires scoring the rivals, and the rivals are usually not published claims that an audit of one paper can reach.

% ------------------------------------------------------------------
\subsection{Interpretive validity}
\label{sec:related-interpretive}

Three of the interpretive criteria have no counterpart we have located in the validity frameworks surveyed above. V1 asks whether a claim declares the description mode at which it is pitched, V2 whether its evidence supports a claim at that level, and V3 whether a simpler description would account for the same evidence. The remaining two are less novel and we do not claim them: scope declaration (V5) is required by every reporting standard in medicine, and unlicensed labeling (V4) is arguably a special case of construct validity.

The absence of V1--V3 elsewhere is explicable. A treatment effect cannot be pitched at the wrong level of description, and a test score either measures its construct or does not. A claim about how a system works can be well measured, causally established, and still wrong about what its components do, because the label attached to a component asserts more than the intervention that identified it can support. Marr's levels and the mechanistic literature on levels of organization supply the conceptual apparatus for describing this---\citet{shagrir2017marr} analyze how Marr's computational level delineates the phenomenon a mechanistic explanation targets---but levels there are a taxonomy of description rather than a criterion an auditor applies to evidence. Turning the taxonomy into a scored criterion is the step we take.

% ------------------------------------------------------------------
\subsection{Completeness and support}
\label{sec:related-completeness}

\citet{craver2006when} states both of the distinctions this dimension rests on. Mechanistic models ``can lie anywhere on a continuum between a mechanism sketch and an ideally complete description of the mechanism,'' where a sketch ``characterizes some parts, activities, or features of the mechanism's organization, but it has gaps'' that are often ``masked by filler terms.'' Separately, how-possibly models are ``only loosely constrained conjectures about the mechanism that produces the explanandum phenomenon'' while how-actually models ``describe real components, activities, and organizational features of the mechanism that in fact produces the phenomenon.''

Whether the second distinction is ontological or epistemic is contested, and the answer matters for what a tier can mean. Craver's own language is ontological. \citet{weiskopf2011models} argues the dimension must be epistemic, since any one of a set of how-possibly models might turn out to be accurate, so their placement ``cannot just be in terms of accuracy''; the dimension represents ``the degree of confirmation of the claim that the model corresponds to the mechanism.'' \citet{colombo2015models} adopt that reading and note that Craver ``does not spell out how we should determine whether a given model is nearer to the how-possibly or how-actually end.'' We adopt the epistemic reading, because a validity framework rates warrant rather than truth, and we flag it as our reading of Craver rather than his stated position.

What we have not found in Craver or in the commentary is the two distinctions crossed. Sketch-to-complete and how-possibly-to-how-actually are presented in separate sections, and no source we have read claims they vary independently or reports a claim's position on both. Reporting them as a pair is what lets the framework distinguish a well-evidenced sketch from a poorly evidenced schema, which are different results that a single label conflates. Two claims in our case studies reach the same tier from opposite ends of this dimension.

% ------------------------------------------------------------------
\subsection{Sources of the criteria}
\label{sec:related-sources}

Every criterion is borrowed from a field that already enforces it. Convergent, discriminant and nomological validity are \citeauthor{cronbach1955construct}'s; reliability, calibration, sensitivity and invariance are psychometric; selection correction is multiple-comparison control; necessity and sufficiency correspond to \citeauthor{craver2007explaining}'s interference and stimulation experiments and offset coupling to the top-down leg of mutual manipulability; double dissociation is neuropsychological; confound control and confounding sensitivity come from causal inference, and \citet{woodward2003making} supplies the interventionist account they operationalize; epistatic interaction and rescue reversibility are genetic; the recording of insufficiency as a graded conclusion follows the USPSTF \citep{uspstf2021}.

Assembling these into one instrument required adapting each of them, because each was developed for a field with an evidential capacity that mechanistic interpretability does not share. Genetics can perform a complementation test \citep{benzer1955fine} because it can construct the double heterozygote; the analogous construction is available for neural networks (C6); epidemiology cannot examine a mechanism's parts; interpretability can read every weight and cannot observe a system it did not train. What the assembly produces is a single instrument under which a claim's construct definition, its measurement properties, its causal warrant, its reach and its interpretation are assessed together and in a stated order, so that a strong result on one dimension stops standing in for the others.

% ------------------------------------------------------------------
\subsection{Evaluation infrastructure in mechanistic interpretability}
\label{sec:related-mi}

Mechanistic Validity is designed to complement, not replace, existing MI evaluation infrastructure. MIB \citep{mueller2025mib} is method-vs-method: which discovery algorithm produces the best circuit? TracR \citep{lindner2024interp} is method-vs-ground-truth: does the method find the true circuit in a synthetic model? Mechanistic Validity is claim-centric: given a claim about how a model works, what evidence supports it and what remains untested? MIB and TracR evaluate discovery methods; Mechanistic Validity evaluates the claims built on their outputs.

Recent work applies construct validity theory to ML evaluation more broadly. \citet{bean2025measuring} find that only 16\% of 445 reviewed LLM benchmarks conduct any statistical testing---a striking parallel to our finding that most MI claims lack measurement validity (M1--M7). Bean et al.'s diagnosis is that benchmarks are treated as self-validating once they gain community adoption; the same dynamic applies to MI metrics, where faithfulness and IIA are used as primary evidence without calibration checks. \citet{freiesleben2025benchmarking} develop an argument-based construct-validity framework for benchmark evaluation, following Messick and Kane, arguing that benchmarks should be evaluated as measurement instruments with explicit reliability and validity properties. Mechanistic Validity differs in two respects: the unit of analysis is a mechanistic claim (a causal assertion about how a model works, not a task-performance score), and the validity structure extends from construct validity alone to five types in dependency order---construct, measurement, internal, external, and interpretive. Where Freiesleben et al.\ ask ``does this benchmark measure the right thing?'' and Bean et al.\ ask ``is this benchmark statistically sound?,'' Mechanistic Validity asks ``does the evidence pattern support the causal claim?''

SAEBench \citep{karvonen2025saebench} audits SAE evaluation metrics; Mechanistic Validity treats that audit as evidence about its own measurement-validity criteria (M1--M7). MechEvalAgent \citep{bai2026mecheval} checks whether published MI code reproduces; Mechanistic Validity asks whether the evidence pattern supports the claim. Unstructured LLM-as-judge evaluation \citep{zheng2023judging} lacks the operational criteria to distinguish validity from plausibility. A companion paper automates Mechanistic Validity via structured LLM extraction against the 36 criteria. \citet{mahale2026causally} likewise grounds interpretability in causal analysis, but its target is the \emph{faithfulness of a natural-language explanation} of an already-discovered circuit, scored with ERASER-style sufficiency and comprehensiveness; Mechanistic Validity instead asks whether the evidence pattern supports the causal \emph{claim} itself, decomposed across all five validity types rather than the single axis of explanation faithfulness.

\citet{sharkey2025open} survey the field's open problems and count the validation of interpretability findings among them. Mechanistic Validity supplies an instrument for one part of that problem: what a given body of evidence establishes about a single claim. \citet{orgad2026actionable} propose a different standard, judging interpretability by \emph{actionability}---whether a finding enables concrete decisions beyond interpretability research itself. The two standards come apart, and the gap between them is this paper's subject. Actionability asks whether a finding is usable; validity asks whether the evidence warrants believing it. A finding can be actionable and unwarranted, and the criteria developed here are what distinguish those cases.

The Schmidt Sciences Trustworthy AI agenda \citep{schmidt2026trustworthy} calls for construct validity, predictive validity, and cross-model generalization of interpretability findings---using these terms but without a unified framework for operationalizing them. Mechanistic Validity provides the operationalization: each of those desiderata maps onto specific criteria (C1--C6 for construct validity, E3--E4 for predictive validity, E4 for cross-model generalization) with concrete metrics and tier requirements.

% ============================================================
\section{Example: From Claim to Verdict}
\label{sec:overview}

Consider a concrete claim: ``The IOI circuit uses three name-mover heads (9.9, 9.6, 10.0) to copy the indirect object to the output position.'' This claim appears in a widely-cited paper, is supported by ablation experiments, and has a clear mechanistic narrative. What has this claim established, and what remains untested?

The framework evaluates such a claim through seven layers:

\begin{enumerate}[nosep]
\item \textbf{Description mode.} At what level is the claim made? The IOI claim operates at the implementational-functional level: it names specific heads and says what each one does.

\item \textbf{Evidence family.} What kinds of evidence are available? The original IOI paper provides primarily interventional-activation evidence (ablation, patching) and some behavioral evidence (faithfulness). Weight-based and cross-model evidence would strengthen the claim by adding independent failure modes.

\item \textbf{Metrics.} What concrete measurements were taken? Activation patching, role ablation, and faithfulness scores. Each metric produces a number; the question is whether those numbers are calibrated and comparable across methods (\Cref{sec:methodology}).

\item \textbf{Criteria.} Does the evidence meet the bar? Sufficiency holds on average (I2: mean-ablating everything outside the circuit leaves 87\% of the logit difference), while necessity at the component level is weaker than the narrative suggests---knocking out all three name-mover heads costs a 5\% drop, because backup heads take over. Specificity (I4: does the circuit affect \emph{only} IOI?) has been tested, with answers that disagree across models. And the faithfulness number is method-conditional in both directions: the same quantity spans below 0\% to over 100\% across six methodological choices \citep{miller2024transformer}, which leaves intervention reach (E1) inconclusive rather than partially satisfied.

\item \textbf{Validity type.} Which validity dimensions are satisfied? Construct validity is partial (the claim is well-defined; discriminant validity is contested). Measurement validity fails on two criteria---stability and invariance both have negative results. Internal validity is bounded by that failure, and external validity has a negative result of its own on prompt generalization.

\item \textbf{Synthesis.} How is evidence aggregated across methods? Currently it is not: separate discovery methods return overlapping but non-identical head sets, and no procedure adjudicates between them. Applying cross-method synthesis (e.g., Dawid-Skene consensus across patching, weight analysis, and information-theoretic methods) would determine which heads survive reliability-weighted voting.

\item \textbf{Verdict.} What tier does the claim reach? Causally Suggestive---the localization is real and reproducible, and the claim is scored at its most lenient reading, on which rival head sets reaching comparable faithfulness leave it standing. Disagreement across ablation methods (E1) is what caps it.
\end{enumerate}

The 36 criteria also formalize common failure modes observed in published MI research (\Cref{tab:failures}).

\begingroup
\small
\begin{table}[H]
\centering
\caption{Ten failure modes observed in published MI research, each mapped to the criteria that detect it.}
\label{tab:failures}
\renewcommand{\arraystretch}{1.15}
\begin{tabular}{@{}l>{\raggedright\arraybackslash}p{3.2cm}>{\raggedright\arraybackslash}p{2.8cm}>{\raggedright\arraybackslash}p{4.0cm}@{}}
\toprule
\textbf{Failure mode} & \textbf{What happens} & \textbf{Criterion} & \textbf{Example} \\
\midrule
Cherry-picking metrics & Run multiple ablation methods, report only the best number & E1 Intervention\newline reach & IOI: one ablation value at origin; the same quantity spans below 0\% to over 100\% across methods \citep{miller2024transformer} \\
\midrule
No baseline control & Report high score without testing a random or untrained model & M2 Baseline\newline separation & SAEBench: Targeted Probe Perturbation scores \emph{higher} on a randomly initialised model than on the trained one \citep{karvonen2025saebench} \\
\midrule
No negative controls & Only test what should pass; never test what should \emph{not} be affected & I4 Specificity\newline C1 Falsifiability & Docstring circuit: effect on non-docstring behavior never measured, in a model that also handles natural language \\
\midrule
Untested discriminant validity & Separation from a neighboring construct is presupposed, never measured & C4 Discriminant\newline validity & Gender bias circuits: gender competence is never localized, so no overlap statistic exists \\
\midrule
Off-target effects reported but not read & Off-target numbers are measured, printed, and left uninterpreted & I4 Specificity & Knowledge neurons: fact updating raises inter-relation perplexity 7.2 against 4.3 for random neurons \citep{dai2022knowledge} \\
\midrule
\midrule
Post-hoc role labeling & Name heads after observed behavior; the label cannot fail & C1 Falsifiability & ``Name mover'' assigned after seeing logit attribution \\
\midrule
No double dissociation & Show necessity but never test the converse & I6 Double\newline dissociation & IOI circuit ablated $\to$ IOI breaks; SVA circuit $\to$ IOI intact never tested \\
\midrule
Interpretive inflation & Label implies intent or algorithm beyond what evidence shows & V4 Unlicensed\newline labeling & ``World model'' for Othello, never contrastively tested against a causally-used state summary \\
\midrule
Scope creep & Claim reach beyond the systems tested, with no experiment behind it & V5 Scope decl.\newline E4 Cross-model & Knowledge neurons: ``not limited to BERT and can be easily generalized to other pretrained models,'' with no second model run \\
\midrule
\midrule
Ignoring alternatives & Circuit works, but so does a different set of heads & I5 Rival mech.\newline exclusion & \citet{chen2026circuits} exhibit two IOI subgraphs at 100\% accuracy sharing 4.1\% of their edges \\
\bottomrule
\end{tabular}
\end{table}
\endgroup


% ============================================================
\section{The Framework}
\label{sec:framework}

The Mechanistic Validity framework evaluates mechanistic claims through seven layers in pipeline order. Each layer answers a distinct question: what kind of explanation is offered (description mode), what sources of signal support it (evidence families), what was concretely measured (metrics), whether the evidence meets specific conditions (criteria), which dimensions of validity are addressed (validity types), how evidence is aggregated (synthesis), and what the claim has established (verdict). 

The framework defines the structure of evaluation---what questions must be answered and in what order. The methodology (\Cref{sec:methodology}) defines the concrete procedures for generating evidence: 59 metrics, including 15 calibrations. The framework evaluates; the methodology produces what is evaluated.

The framework's layers are described in the subsections that follow: description modes (\Cref{sec:description-modes}), evidence families (\Cref{sec:evidence-families}), metrics (\Cref{sec:metrics}), criteria (\Cref{sec:criteria}), validity types (\Cref{sec:validity-types}), synthesis protocols, and verdict tiers (\Cref{sec:verdict-tiers}).


\subsection{Description modes}
\label{sec:description-modes}

A description mode is a commitment about what kind of explanation is being offered (\Cref{tab:modes}). The declaration comes before the analysis and determines what counts as evidence, what counts as a gap, and what the finding means if it succeeds. The distinction matters because mechanistic interpretability routinely conflates levels: a paper that discovers \emph{which components} are active (implementational-topographic) writes a narrative about \emph{what the model computes} (computational). Each level upgrade requires additional bridging evidence that the lower level cannot provide. Anthropic's global-workspace result is the sharpest recent instance, and recurs below as a running example: the Jacobian lens recovers a residual-stream subspace (implementational-statistical), the paper says that subspace encodes what the model can report (representational), and it narrates the whole as a ``global workspace'' the model reasons in (computational)---three modes deep, on evidence that reaches the representational level at best \citep{gurnee2026verbalizable}. Declaring the mode first forces the question the label skips: what bridging evidence carries a subspace claim up to a computational one?

The seven modes extend Marr's three-level hierarchy \citep{marr1982vision} by splitting the implementational level into four sub-types---functional, connectomic, statistical, and topographic---reflecting the fact that MI papers routinely conflate claims at these different levels. The modes form a partial order by evidential commitment. A computational claim commits to a function being computed and a reason it is the right function. An algorithmic claim commits to a specific procedure. A representational claim commits to what information is encoded but not how it is processed. The four implementational sub-types commit to progressively weaker structural facts: which components are involved (topographic), how they are wired (connectomic), what their activations look like (statistical), and what input-output transformation each performs (functional).

\begin{table}[H]
\centering
\small
\caption{Seven description modes, ordered from most committal (computational) to least committal (topographic). Higher modes require all evidence from lower modes plus bridging evidence. Implementational-statistical is partially orthogonal: it characterizes activation distributions rather than making structural commitments.}
\label{tab:modes}
\renewcommand{\arraystretch}{1.2}
\begin{tabular}{@{}rl>{\raggedright\arraybackslash}p{4.8cm}>{\raggedright\arraybackslash}p{5.2cm}@{}}
\toprule
& \textbf{Mode} & \textbf{What it asks} & \textbf{Example} \\
\midrule
1 & Computational & What function does the system compute, and why? & ``GPT-2 predicts the indirect object'' \\
\midrule
2 & Algorithmic & What procedure does it execute? & ``A detect-inhibit-copy algorithm'' \\
\midrule
3 & Representational & What information is encoded, and in what geometry? & ``The IO name is linearly encoded at layer 8'' \\
\midrule
4 & Impl.-Functional & What transformation does each component perform? & ``Head 9.9 copies names to the output'' \\
\midrule
5 & Impl.-Connectomic & How are components wired? & ``S-inhibition feeds name movers via residual stream'' \\
\midrule
6 & Impl.-Statistical & What do activations look like? & ``Head 9.9 attends strongly to IO position'' \\
\midrule
7 & Impl.-Topographic & Which components are involved? & ``Heads 9.9, 9.6, 10.0'' \\
\bottomrule
\end{tabular}
\end{table}

A claim may carry evidence at several modes simultaneously; the framework does not average them. Each mode-level sub-claim receives its own assessment, and the overall verdict is reported at the strongest mode whose required evidence is complete, with lower-mode evidence noted as supporting and higher-mode evidence as suggestive. The mode assigned to a verdict is thus the strongest level for which all required evidence is present. Partial evidence at a higher level does not upgrade the mode---it is reported as suggestive while the verdict is issued at the established mode. For example, a claim at the algorithmic level (``the model runs a detect-inhibit-copy algorithm'') requires both the topographic evidence (which heads are involved) and the additional evidence that the heads implement the claimed procedure---path-level causal tracing, timing consistency, and sufficiency of the minimal circuit.

\subsection{Evidence families}
\label{sec:evidence-families}

An evidence family classifies a metric's output by the source of signal it draws on, independent of the specific tool used to generate it. Four families exhaust the sources available for a mechanism claim: \textbf{structure}, the system's persistent organization, observable without running it; \textbf{state}, what the system does internally while operating; \textbf{behavior}, its input-output relation; and \textbf{history}, how it came to be as it is. In mechanistic interpretability these are weights, activations, task performance, and training, and the rest of this paper uses those names, since the case studies are drawn from that field. Each family crosses with a second axis---observational vs.\ interventional---producing eight cells (\Cref{tab:families}). Most published MI evidence lives in a single cell (interventional activations: patching and ablation). Triangulation means covering multiple cells, and the observational/interventional axis matters as much as the family axis for independence of failure modes: two interventional-activation metrics share confounds that an observational-weight metric does not.

\begin{table}[H]
\centering
\small
\caption{Evidence families $\times$ evidence mode. Each cell contains representative methods. Convergent evidence across cells with independent failure modes is qualitatively stronger than redundant evidence from the same cell. Most published MI work concentrates in interventional activations (bold).}
\label{tab:families}
\renewcommand{\arraystretch}{1.25}
\begin{tabular}{@{}l>{\raggedright\arraybackslash}p{5.8cm}>{\raggedright\arraybackslash}p{5.8cm}@{}}
\toprule
& \textbf{Observational} & \textbf{Interventional} \\
\midrule
\textbf{Weights} & SVD, effective rank, OV/QK composition, minimum description length & Circuit transplant, weight editing, weight knockout \\
\midrule
\textbf{Activations} & Probing, CKA, mutual information, partial information decomposition, logit lens & \textbf{Activation patching, DAS-IIA, ablation, steering} \\
\midrule
\textbf{Behavior} & Logit diff, behavioral profiling, task accuracy & Prompt perturbation, dose-response, KL under ablation \\
\midrule
\textbf{Training} & Loss curves, checkpoint comparison, phase transitions & Ablate-then-retrain, fine-tuning, curriculum manipulation \\
\bottomrule
\end{tabular}
\end{table}

Each family answers a different question about the mechanism and has characteristic strengths and blind spots. Weight-based metrics cannot be confounded by runtime compensation but cannot distinguish used capacity from unused capacity. Activation-based metrics can establish necessity and sufficiency directly but are vulnerable to backup circuits that compensate after ablation. Behavioral metrics test whether the circuit reproduces the model's outputs, but behavioral equivalence does not imply mechanistic equivalence---different circuits can produce the same behavior for different reasons. Training-based metrics reveal developmental trajectories but cannot establish that a mechanism is causally responsible for current behavior.

The four families hold their shape across fields, and what each field can put in them differs (\Cref{tab:families-domains}). The empty cell is as informative as the full ones: epidemiology has no structural evidence in this sense, because there is no persistent organization to inspect apart from the system running, which is why the field leans on history and state and built its deepest machinery there. A field that can access only three families is capped on convergent validity (C3) by construction rather than by neglect.

\begin{table}[H]
\centering\small
\caption{The same four evidence families across fields. Entries name what each field can measure in that family; a dash marks a family the field cannot access.}
\label{tab:families-domains}
\renewcommand{\arraystretch}{1.2}
\begin{tabular}{@{}l>{\raggedright\arraybackslash}p{2.2cm}>{\raggedright\arraybackslash}p{2.2cm}>{\raggedright\arraybackslash}p{2.2cm}>{\raggedright\arraybackslash}p{2.2cm}>{\raggedright\arraybackslash}p{2.4cm}@{}}
\toprule
& \textbf{Interpretability} & \textbf{Genetics} & \textbf{Neuroscience} & \textbf{Epidemiology} & \textbf{Medical microbiology} \\
\midrule
\textbf{Structure} & Weights & Genome, protein structure & Connectome, anatomy & --- & Pathogen sequence \\
\midrule
\textbf{State} & Activations & Expression levels & Firing rates, fMRI & Biomarkers & Copy number, load \\
\midrule
\textbf{Behavior} & Task performance & Phenotype & Behavioral task & Outcome, disease status & Clinical disease \\
\midrule
\textbf{History} & Training checkpoints & Development, lineage & Maturation, plasticity & Exposure history & Infection course, relapse \\
\bottomrule
\end{tabular}
\end{table}

Some methods straddle multiple families. Gradient-based attribution (EAP, integrated gradients) computes counterfactual signals through both weights and activations jointly---it is neither purely observational nor interventional. Causal scrubbing intervenes on activations but evaluates through behavioral output. Automated circuit discovery methods like ACDC compose interventional-activation and behavioral evidence across multiple rounds. The family classification assigns each method to its \emph{primary} source of signal; methods that draw on multiple families are noted as such, and their cross-family nature is itself informative for convergent validity assessments.

Within each family, interventional evidence is epistemically stronger than observational evidence because it rules out confounds that observation alone cannot. When evaluating a claim, the analyst asks: how many cells are covered? Which families and modes are represented? And are there cells that \emph{should} support the claim but do not? Absence of expected evidence is informative---it may indicate a gap in the claim or a limitation of the proposed mechanism. Measurement calibrations---bootstrap stability, seed variance, baseline separation---cut across all eight cells and are described in \Cref{sec:calibrations}.

\subsection{Criteria}
\label{sec:criteria}

Criteria are specific, falsifiable conditions that must be met for a validity type to be satisfied (\Cref{tab:criteria}). Each criterion has a concrete pass condition---for example, I1 (necessity) requires that ablating the component degrades performance across at least two independent methods, not just one. A claim either passes, partially passes, or fails each criterion, producing a structured validity profile rather than a single score.

\begingroup
\hyphenpenalty=10000
\exhyphenpenalty=10000
\begin{table}[H]
\centering
\small
\caption{Thirty-six criteria across five validity types. Each criterion has a concrete test; a claim either passes, partially passes, or fails, producing a structured validity profile.}
\label{tab:criteria}
\renewcommand{\arraystretch}{1.15}
\begin{tabular}{@{}llp{8.5cm}@{}}
\toprule
\textbf{ID} & \textbf{Criterion} & \textbf{Description} \\
\midrule
C1 & Falsifiability & Can the claim be refuted? \\
C2 & Structural plausibility & Is the mechanism physically possible in the architecture? \\
C3 & Convergent validity & Do multiple independent methods agree? \\
C4 & Discriminant validity & Does the measure distinguish this from neighboring constructs? \\
C5 & Nomological validity & Does the claim fit into a broader theory? \\
C6 & Complementation validity & Are the construct's labeled subdivisions functionally distinct? \\
\midrule
M1 & Reliability & Do repeated measurements give the same answer? \\
M2 & Baseline separation & Is the score distinguishable from random/untrained baselines? \\
M3 & Stability & Is the classification robust to perturbation? \\
M4 & Calibration & Are the numbers meaningful? \\
M5 & Sensitivity & Can the instrument detect known-true effects? \\
M6 & Invariance & Does the metric behave consistently across conditions? \\
M7 & Selection correction & When $k$ findings are selected from $N$ candidates, is $N$ reported and multiplicity controlled? \\
\midrule
I1 & Necessity & Is the circuit required for the behavior? \\
I2 & Sufficiency & Is the circuit enough to produce the behavior? \\
I3 & Minimality & Does every component earn its place? \\
I4 & Specificity & Does intervening on the circuit affect this task more than matched control tasks? \\
I5 & Rival mechanism exclusion & Is this \emph{the} mechanism, or \emph{a} mechanism? \\
I6 & Double dissociation & Do two interventions cross, each breaking what the other spares? \\
I7 & Confound control & Are alternative explanations ruled out? \\
I8 & Confounding sensitivity & How strong must an unmeasured confounder be to explain the result? \\
I9 & Epistatic interaction & Do circuit components interact non-additively, and does the direction of interaction distinguish shared pathways from mutual compensation? \\
I10 & Rescue reversibility & Does restoring a corrupted component recover behavior? \\
I11 & Onset coupling & Does the mechanism appear when the capability appears? \\
\midrule
I12 & Offset coupling & Does the mechanism go when the capability is removed? \\
\midrule
E1 & Intervention reach & Has the result been reproduced under at least two intervention families, and do they agree? \\
E2 & Prompt generalization & Does it work on diverse prompts? \\
E3 & Cross-task generalization & Does the mechanism transfer to related tasks? \\
E4 & Cross-model generalization & Does the mechanism appear in other models? \\
E5 & Graded response & Does partial ablation produce partial effects? \\
E6 & Novel prediction & Does the mechanism predict new, untested behaviors? \\
\midrule
V1 & Level declaration & At what description mode is the claim made? \\
V2 & Level-evidence match & Does the evidence support claims at that level? \\
V3 & Alternative level & Could the evidence be explained at a different level? \\
V4 & Unlicensed labeling & Does a name import a property that was not measured? \\
V5 & Scope declaration & What does the claim explicitly not cover? \\
\bottomrule
\end{tabular}
\end{table}
\endgroup

Two of the internal-validity criteria address what necessity and sufficiency leave open. Both are properties of the component set as a whole, so a set can satisfy both while containing members that do nothing: dropping a redundant head changes neither the ablation deficit nor the recovered logit difference. Minimality (I3) is the test that catches this. Sufficiency is likewise satisfied by more than one set---\citet{chen2026circuits} exhibit two IOI circuits at perfect accuracy sharing 4.1\% of their edges---so a passing intervention shows the reported circuit to be \emph{a} sufficient mechanism. Rival mechanism exclusion (I5) asks for the stronger result, that no comparable rival survives the same test. The block is ordered by what each criterion asks. I1--I3 are properties of the set as a whole; I4--I6 discriminate across tasks, across rival circuits, and across both; I7--I8 cover measured and unmeasured confounders; I9--I10 probe internal structure.

The internal-validity criteria are written in terms of interventions because interventions are the strongest evidence for necessity and sufficiency. Fields that cannot intervene establish I1 and I2 by other means: natural experiments, instrumental variables, comparison across cases. Those routes reach the same conclusion at the cost of assumptions the intervention did not need, so a claim should record which route it took.

V5 records the scope a claim asserts. Whether the evidence reaches that scope is a separate question, addressed in the companion work on admissibility.




\begingroup
\hyphenpenalty=10000
\exhyphenpenalty=10000
\begin{table}[H]
\centering
\small
\caption{Criterion status categories. Status labels summarize the evidential state of an individual criterion, not the overall verdict for the claim.}
\label{tab:criterion-status}
\renewcommand{\arraystretch}{1.15}
\begin{tabular}{@{}lp{10.5cm}@{}}
\toprule
\textbf{Status} & \textbf{Meaning} \\
\midrule
Confirmed & Tests capable of exposing the relevant failure mode were conducted and passed. \\
\midrule
Partially confirmed  & Supportive evidence exists, but it is too narrow, method-dependent, weakly calibrated, or insufficiently severe to establish the criterion. \\
\midrule
Inconclusive & Relevant tests give conflicting results across methods, datasets, interventions, calibrations, or scopes. \\
\midrule
Disconfirmed & An appropriate test was conducted and the criterion failed. \\
\midrule
Untested & No test capable of evaluating the criterion has been conducted. \\
\bottomrule
\end{tabular}
\end{table}
\endgroup

For each validity criterion we assign one of six statuses: Confirmed, Partially confirmed, Inconclusive, Untested, Disconfirmed, Not applicable. Confirmed and Partially confirmed establish a criterion. The other three do not, and tier promotion sees no difference among them. The remaining three are not ranked against one another: a criterion that was tested and failed and one that nobody tested both leave the criterion unestablished, and ordering them would make not running a test preferable to running one and reporting an adverse result. These statuses are ordinal verdict categories in the error-statistical tradition of \citet{mayo1996error, mayo2018statistical}: a criterion is Confirmed if and only if the relevant experimental tests were conducted under conditions of sufficient severity to have detected failure, and the criterion passed those tests. Following \citet{achinstein2001book}, we further distinguish criteria where the available evidence constitutes a genuine epistemic reason to accept the criterion (Confirmed) from those where evidence is present but falls short of that bar (Partially confirmed, Inconclusive).

This framework is non-probabilistic: status assignments do not represent posterior beliefs or confidence levels. Instead, following \citet{campbell1963experimental} and \citet{cook1979quasi}, validity claims are evaluated by whether the relevant tests rule out plausible rival explanations. The Untested status applies when no test capable of evaluating the criterion has been conducted, whereas Inconclusive indicates contested evidence and Disconfirmed indicates that a test was conducted and the criterion failed. Not applicable is reserved for criteria the claim’s structure cannot pose: a toy model trained on one task has no off-target behavior against which to measure specificity, and a released checkpoint has no training history against which to measure onset. It is argued from what the claim is, never from what an experiment would cost, and it is excluded from the denominator rather than counted as passing.

Recording insufficiency as a verdict has an institutional precedent. The US Preventive Services Task Force issues an \emph{I statement} when the evidence is insufficient to assess the balance of benefits and harms, making inadequacy a standard graded conclusion rather than a remark \citep{uspstf2021}. Null and negative results are of course published in mechanistic interpretability; what the Untested and Inconclusive statuses add is a place for them in the same ordinal scale as positive findings, so that ``this was never tested'' and ``the tests disagree'' are recorded in the audit rather than left to prose. Our scale divides the USPSTF's single category: an I statement covers both a question never tested and one whose tests conflict, which we separate as Untested and Inconclusive. The separation earns its cost because the remedies differ. An Untested criterion needs an experiment; an Inconclusive one needs adjudication between results that already exist.


These statuses evaluate individual criteria, not the overall claim. A claim can therefore contain Confirmed internal-validity criteria, Untested measurement-validity criteria, and Mixed external-validity criteria while still receiving a single overall verdict tier. This keeps criterion-level strengths, gaps, and failures explicit instead of compressing them into an undifferentiated verdict.

\subsection{Validity types}
\label{sec:validity-types}

The five validity types form a dependency chain: later types cannot be established without first satisfying earlier ones. The ordering reflects logical precedence: a construct that is not well-defined cannot be reliably measured, an unreliable measurement cannot support a causal claim, a causal claim in one setting cannot be generalized, and an interpretation cannot be evaluated until the underlying evidence is characterized. The vocabulary is \citeauthor{shadish2002experimental}'s, the ordering is not. Their organization pairs statistical conclusion validity with internal validity and construct validity with external validity, placing construct third, and does not present the four as a strict precedence relation. We place construct first, following \citet{messick1995validity}, for whom construct validity ``undergirds all score-based interpretations'' and is ``the essence of a unitary validity conception.'' A claim whose target concept is incoherent cannot be rescued by any amount of causal evidence, so the construct has to come first for the same reason Messick gives. Reading the remaining four as a strict chain is our own commitment, and the audits are what test it: if a claim could establish external validity while failing measurement validity, the chain would be wrong.

One qualification on the second link. Reliability (M1) bounds internal validity in a strict sense: classical test theory's attenuation result caps the correlation two measures can exhibit at the geometric mean of their reliabilities, so an unreliable instrument limits how much causal validity is demonstrable at all. Stability (M3) and invariance (M6) do not bound in that sense. They record that a reported quantity moves with analysis choices or across conditions, which impugns the number that was published without ceilinging the effect it estimates. A circuit can be genuinely necessary while the faithfulness figure reported for it is one of several the protocol could have produced. We therefore let M1 cap and report M3 and M6 as qualifications on the evidence.

One criterion is scored differently from the rest, and the reason generalizes. A double
dissociation (I6) is a single crossed design rather than an accumulation of evidence: two
mechanisms, two behaviors, each intervention moving one and sparing the other. Half of it
earns no partial credit, because each half is already scored elsewhere --- the arm in which
disrupting the mechanism breaks the behavior is necessity (I1), and the arm in which a
control intervention spares it is baseline separation (M2). Awarding a partial pass at I6
for either would count one experiment twice and promote a claim on evidence the ladder has
already used. We therefore score I6 as met or unmet. Across the sixteen audited claims it
is met once, by a test run three years after the origin paper and by different authors
\citep{feucht2025dual}, which is the sharpest single measurement of how rarely the field
runs the design its strongest tier requires.

\begin{equation}
\text{Construct} \rightarrow \text{Measurement} \rightarrow \text{Internal} \rightarrow \text{External} \rightarrow \text{Interpretive}
\label{eq:dependency}
\end{equation}

The pipeline---description mode, evidence family, metrics, criteria, validity type---is procedural---it describes how evidence is built or audited. The validity-type chain is logical---it describes which types presuppose which. At pipeline step 5 (validity type), the framework evaluates every validity type the claim asserts, but the dependency order constrains the verdict: a claim cannot be credited with external validity beyond the level its measurement and internal validity support. The chain is a ceiling, not a sequence of gates.

Each type addresses a distinct question about the claim:

\paragraph{Construct validity.} Is the thing being measured well-defined? Before any measurement is taken, the construct must be specified precisely enough that the claim is falsifiable, structurally plausible, and distinguishable from neighboring constructs. A ``deception feature'' that cannot be distinguished from an ``uncertainty feature'' fails construct validity regardless of how carefully it is measured. This draws on philosophy of science (Popper's falsifiability, Bridgman's operationalism) and psychometrics (construct validity, convergent and discriminant validity).

\paragraph{Measurement validity.} Are the instruments trustworthy? \citet{borsboom2004concept} argue that a measure is valid if and only if the attribute exists and variation in it causally produces variation in the measurement outcome. This causal account grounds the framework's measurement criteria. Once the construct is defined, the measurement tools used to detect it must be reliable (stable across repetitions), calibrated (separable from random baselines), and invariant (consistent across conditions). A metric that gives different answers when run with different random seeds is not evidence. This draws on psychometrics (test-retest reliability, measurement invariance) and pharmacology (assay validation).

\paragraph{Internal validity.} Does the evidence support the causal claim? Given a well-defined construct and trustworthy instruments, the evidence must establish that the identified component is causally involved in the behavior---not merely correlated with it. This requires demonstrating necessity (the component is required), sufficiency (the component is enough), and specificity (the component does this task and not everything). Critically, a component that is causally necessary for many tasks is a bottleneck, not an implementation---double dissociation, not single dissociation, is the test for specificity. This draws on neuroscience (lesion vs.\ stimulation, single vs.\ double dissociation, mutual manipulability; \citealp{shallice1988neuropsychology, craver2007explaining}), causal inference (do-calculus, counterfactual frameworks, instrumental variables; \citealp{pearl2009causality}), and genetics (knockout vs.\ rescue, epistasis mapping for non-additive component interactions, sensitivity analysis for unmeasured confounders).

\paragraph{External validity.} Does the mechanism generalize? Causal evidence on one prompt set, with one ablation method, in one model, establishes internal validity at best. External validity requires that the mechanism generalizes across prompts, ablation methods, related tasks, and ideally across models. Beyond replication, it requires characterizing the effect quantitatively: how it scales with intervention strength (the dose-response curve), whether it is selective for the claimed task versus related tasks (the selectivity ratio), and whether its absolute magnitude supports the computational story. This draws on pharmacology (Phase~III generalization, dose-response, therapeutic windows; \citealp{rang2006drug}) and causal inference (transportability; \citealp{pearl2011transportability}).

\paragraph{Interpretive validity.} Is the interpretation of the mechanism correct? The natural-language description of the mechanism must be justified by the evidence. A head called a ``name mover'' based on direct logit attribution carries theoretical implications beyond what the evidence strictly supports. Interpretive validity requires declaring the Marr level at which the claim operates---implementational, algorithmic, or computational---matching evidence to that level, and checking that the narrative is consistent across evidence families. \citet{meloux2025everything} prove that on Boolean MLPs small enough to enumerate, multiple circuits replicate the same behavior and multiple interpretations fit the same circuit; alternative exclusion is therefore a required step, not an optional one. This draws on Messick's unified theory of validity \citep{messick1995validity}---which frames validity as the warrant for \emph{interpretations} of scores, not properties of the test itself---and on Marr's levels of analysis \citep{marr1982vision}.

\subsection{Verdict tiers}
\label{sec:verdict-tiers}

Claims are assigned to one of five verdict tiers representing qualitative transitions in evidential status (\Cref{tab:verdicts}). The tiers implement the epistemic reading of \citeauthor{craver2006when}'s how-possibly to how-actually dimension discussed in \Cref{sec:related-completeness}: Proposed corresponds to a loosely constrained conjecture and the upper tiers to a model well confirmed by the available evidence. Following \citet{colombo2015models}, a claim must also be the best confirmed among the available accounts of its phenomenon, which the tiers do not measure and the Underdetermined label records. We report the tier together with the criteria that cap it, as \citet{guyatt2008grade} report a certainty rating alongside the domains that lowered it.

\begin{table}[H]
\centering
\small
\caption{Verdict tiers with minimum evidence requirements. Each tier subsumes all requirements of lower tiers. Underdetermined, Insufficient and Disconfirmed replace the tier rather than accompanying it: a claim carrying one of them does not also hold a position in the hierarchy.}
\label{tab:verdicts}
\label{tab:tier-promotion}
\renewcommand{\arraystretch}{1.15}
\begin{tabular}{@{}>{\raggedright\arraybackslash}p{2.2cm}>{\raggedright\arraybackslash}p{3.4cm}>{\raggedright\arraybackslash}p{7.0cm}@{}}
\toprule
\textbf{Tier} & \textbf{Meaning} & \textbf{Minimum evidence requirements} \\
\midrule
Proposed & Structural or representational evidence only & Construct defined under a declared description mode (C1--C2); at least one admissible measurement conducted \\
\midrule
Causally Suggestive & Necessity shown, sufficiency not established & Necessity via causal intervention (I1 confirmed); at least one measurement passes baseline separation (M2) \\
\midrule
Mechanistically Supported & Necessity + sufficiency with consistent methods & Sufficiency established (I2); intervention reach across $\geq 2$ ablation methods (E1); specificity test conducted (I4 at least partially confirmed) \\
\midrule
Triangulated & Multiple converging lines of independent evidence & Convergent evidence from $\geq 3$ evidence families (C3), where families count as independent when the failure of one's core assumption would not automatically invalidate the other; cross-distribution replication---E4 where the claim asserts reach beyond the systems tested, E2 where it does not; double dissociation attempted (I6). Whichever criterion promotes, the status of the other is reported \\
\midrule
Validated & All five validity types addressed & Measurement calibration (M1--M6) explicitly addressed; interpretive validity (V1--V5) audited; external validity across models (E4) and prompts (E2) \\
\midrule
Underdetermined & Evidence consistent with multiple mechanisms & Cannot resolve between rival specifications \\
\midrule
Insufficient & Cannot be assessed & The construct is not defined well enough, or no admissible measurement exists, to score the claim \\
\midrule
Disconfirmed & Fails decisively on a key criterion & Negative result on a required criterion \\
\bottomrule
\end{tabular}
\end{table}

Every tier carries a specific epistemic meaning: a claim at \emph{Proposed} has its evidential status, and the path to strengthening it, precisely characterized. The verdict tells researchers exactly what additional evidence would advance the claim: a Proposed claim needs causal intervention (I1) to reach Causally Suggestive; a Causally Suggestive claim needs sufficiency evidence (I2) and method consistency (E1) to reach Mechanistically Supported.

Some criteria cannot be tested for some claims, and the tier system treats that as information rather than as an exemption. Specificity (I4) asks whether a mechanism does this task and not everything, which a model trained on a single task cannot answer: there is no second task to fail on. Onset--offset coupling (I11) asks whether a mechanism appears when the capability appears, which no claim about a released model without training checkpoints can answer. In each case the claim caps at the tier below, and the cap is a statement about the setting rather than a deficiency in the work. A result on a single-task toy model licenses less than one on a model doing many things, and a field that cannot observe a system's history establishes less than one that can. Waiving the requirement so the tier reads better would reverse the framework's purpose; recording the foreclosure is what makes the limit visible.

Three diagnostic labels sit outside the hierarchy, covering the three ways a claim can fail to occupy a tier: too many explanations fit it, too little evidence exists to assess it, or it has been tested and failed. \emph{Insufficient} follows the US Preventive Services Task Force's \emph{I statement} \citep{uspstf2021}, and marks a claim whose construct is not defined well enough to score rather than one that scores badly. Underdetermined and Disconfirmed serve distinct purposes. Underdetermined means the evidence is consistent with multiple competing mechanisms and the framework cannot resolve between them---an excess of explanations, requiring additional experiments that discriminate between rivals; no claim in this audit carries the label, and the IOI circuit is the instructive near-miss: scoring it at its most lenient reading removes the competition, because rival subgraphs and the audited claim can both hold. Disconfirmed means a criterion has been decisively failed by a test that could have gone the other way---for example, when the construct turns out not to separate from its neighbor (C4), as in the knowledge-neuron case study, where the same editing machinery was later shown to move non-factual linguistic phenomena. Neither label is a statement that no evidence exists; both require evidence, and the difference between them is whether it points in one direction.

\paragraph{Completeness is reported separately from tier.} A verdict states how well a claim is evidenced. It does not state how much of the mechanism the claim specifies, and the two vary independently. The two also trade against each other inside a single paper: every additional component a claim specifies is another component whose role must be evidenced, so mechanistic specificity is bought with evidential risk. We report both. For a fixed body of evidence the achievable combinations form a Pareto frontier---no claim on it can be made more specific without losing tier, or better supported without becoming vaguer---and a paper's contribution is partly which point on that frontier it chose. A claim's position on Craver's continuum is read from structural plausibility (C2) and level--evidence match (V2) rather than scored as a separate criterion. The IOI circuit and superposition both reach Mechanistically Supported and sit at opposite ends of it: one names four of seven head classes by role without a weight-level account, the other derives every component in closed form.

% ============================================================
\subsection{Severity: preregistration and corroboration}
\label{sec:severity}

The top tiers require more than accumulated confirmations. A claim tested only on the evidence it was constructed to fit has been confirmed, not corroborated: it passed a test it could not have failed. Following the error-statistical tradition \citep{mayo2018statistical}, we require the criteria that lift a claim to Triangulated or Validated to be established by \emph{severe} tests---tests the claim could have failed but did not. Severity is earned two ways. Preregistration supplies it directly: a specificity or generalization test whose prediction was fixed in advance is more severe than one selected after seeing the result, so a preregistered pass earns Confirmed where a post-hoc pass earns Partially confirmed. Independent corroboration supplies it structurally: convergent evidence across evidence families with different failure modes (C3) is, by construction, evidence the claim was not optimized to produce. The tier ceiling is therefore a consequence of the existing structure rather than an added penalty---Triangulated already requires cross-family convergence, so a claim resting on a single family, however often re-tested, cannot exceed Mechanistically Supported. Preregistration is the cleaner route to severity, not a separate gate: the verdict label stays a single word, and whether a tier was earned by preregistration or by corroboration is stated in the claim's write-up.

% ============================================================
\section{Methodology}
\label{sec:methodology}

The criteria say what a mechanistic claim owes its reader. This section supplies the
instruments that discharge them, and it runs in both directions. Read forward, it is a
design specification: choose the metrics that bear on the criteria a claim will need,
calibrate them before collecting evidence, and record what each measurement establishes
as it is produced. Read backward, the same instruments recover what a finished paper
established.

The framework layers described in \Cref{sec:framework} define what counts as evidence and how it is evaluated. This section provides the concrete measurement procedures---metrics, calibrations, protocols, and synthesis methods---that generate the evidence evaluated against those layers. These are organized by evidence family and by evidential function.

\subsection{Metrics}
\label{sec:metrics}

Metrics are concrete, runnable tests that produce measurements about a neural network---the layer of the hierarchy where empirical contact with the model actually happens. Everything above (evidence families, criteria, validity types, verdicts) depends on what metrics measure and how well they measure it. A metric alone cannot establish a claim; a claim requires evidence from multiple metrics, evaluated against the criteria appropriate to the validity type being asserted.

Each metric takes a model, a task, and optionally an intervention as input, and produces a number as output. Metrics can be composed: one metric's output feeds into another. The framework organizes evidence into four evidence families by \emph{source of signal}; the catalog below groups metrics by \emph{what they measure}, in six groups A--F, so a group may draw on more than one family. The current catalog defines 59 metrics across groups A--F, drawing on interventionist, information-theoretic, structural, and generalization traditions of causation. The full inventory is in \Cref{appendix:metrics}.

\subsection{Calibrations}
\label{sec:calibrations}

Calibrations are meta-metrics: they take another metric's output as input and assess whether that output is stable (bootstrap stability), reproducible (seed variance), or distinguishable from baselines (baseline separation). A high score on a primary metric is not evidence if the corresponding calibration fails. The 15 calibration metrics correspond to the measurement criteria M1--M7, and \Cref{tab:f-to-m} gives the mapping. The full catalog is in \Cref{appendix:metrics}.

\begin{table}[H]
\centering
\small
\caption{Which measurement criterion each calibration metric bears on.  The
correspondence is many-to-one: reliability and stability each take several
calibrations, and selection correction takes one.}
\label{tab:f-to-m}
\renewcommand{\arraystretch}{1.15}
\begin{tabular}{@{}ll@{}}
\toprule
\textbf{Criterion} & \textbf{Calibration metrics} \\
\midrule
M1 Reliability          & F01 Test--Retest, F02 Bootstrap Stability, F03 Seed Variance, \\
                        & F11 Estimator Uncertainty \\
\addlinespace[2pt]
M2 Baseline separation  & F06 Baseline Separation, F07 Negative Controls \\
\addlinespace[2pt]
M3 Stability            & F04 Checkpoint Variance, F05 Prompt Variance, \\
                        & F09 Intervention Robustness, F10 Hyperparam.\ Sensitivity \\
\addlinespace[2pt]
M4 Calibration          & F12 Calibration Curve, F13 Cross-Metric Convergence \\
\addlinespace[2pt]
M5 Sensitivity          & F08 Positive Controls \\
\addlinespace[2pt]
M6 Invariance           & F14 Measurement Invariance \\
\addlinespace[2pt]
M7 Selection correction & F15 Multiple Comparisons \\
\bottomrule
\end{tabular}
\end{table}


\subsection{Recommended practice}
\label{sec:design}
A validity framework changes nothing if its criteria go unrun, and the fields that developed those criteria supply the measured cost of leaving them unrun. Under the full validation battery psychology itself developed, 4\% of the scales examined passed, against 88\% under the modal practice of reporting internal consistency alone \citep{hussey2020hidden}. Ecology reports the same pattern: 23\% of 537 occupancy studies account for imperfect detection, while measured per-survey detection probability usually falls below 0.5 \citep{kellner2014accounting}. So does epidemiology, where researchers who conducted a quantitative bias analysis frequently left their interpretation unchanged after it moved the estimate materially \citep{petersen2021systematic}.

The audits here place interpretability in the same position. Of the 36 criteria applied to the IOI circuit, twelve are Untested or Inconclusive, and most need no artifact the field lacks: seed variance for the head-evaluator scores, a planted circuit to validate the discovery method, a record of how many candidate circuits were examined, and a double dissociation between the head classes. That gap reflects reporting practice rather than a limit of the method, which is what makes the following recommendations tractable.

The case studies converge on five concrete changes to MI research practice that would have the highest impact on claim quality:
\begin{enumerate}[nosep]
\item \emph{Test specificity by default.} Every circuit ablation study should report performance on at least two unrelated tasks alongside the target task. This is cheap, takes one additional forward pass per task, and addresses the single most common gap (I4) across all sixteen case studies.
\item \emph{Report multiple ablation methods.} A faithfulness number from one ablation method is a lower bound on the causal claim, not the claim itself. At minimum, report mean ablation and resample ablation; ideally add zero ablation and activation patching.
\item \emph{Calibrate metrics against baselines.} Report the metric score on a random circuit of the same size, an untrained model, and a shuffled-label baseline. If the delta between the real score and the baseline is small, the metric is not measuring what it claims.
\item \emph{Declare description mode and scope.} State explicitly whether the claim is computational, algorithmic, or implementational, and what input distribution it is validated on. This costs zero experiments and prevents the most common interpretive inflation.
\item \emph{Report non-uniqueness.} If alternative circuits achieve comparable faithfulness, report this as a finding rather than suppressing it. Non-uniqueness is scientific information; a single reported circuit is a selection from a space the reader cannot see.
\end{enumerate}




% ============================================================
\section{Theoretical Foundations}
\label{sec:foundations}

Each discipline contributes the conditions it requires before an inference is accepted, and names the failure mode it guards against (\Cref{tab:foundations-summary}). The load-bearing transfers appear where they do work rather than in a survey---confirmation versus corroboration in the severity rule (\Cref{sec:severity}), ablation versus patching as formally distinct interventions in the internal-validity criteria, dose-response and system reserve in the external-validity criteria, and description versus explanation in the interpretive criteria.

\begin{table}[h]
\centering\small
\caption{The eight foundations, each grounding a validity type.}
\label{tab:foundations-summary}
\renewcommand{\arraystretch}{1.15}
\begin{tabular}{@{}ll>{\raggedright\arraybackslash}p{7.2cm}@{}}
\toprule
\textbf{Discipline} & \textbf{Grounds} & \textbf{Principal transfer to mechanistic interpretability} \\
\midrule
Philosophy of science & Construct (C) & Falsifiability, severe testing, and the confirmation/corroboration distinction separating a test a claim was built to pass from one it could have failed \\
\midrule
Psychometrics & Construct (C), Measurement (M) & Convergent and discriminant validity from the multitrait-multimethod matrix (C3, C4); reliability and baseline separation \\
\midrule
Causal inference & Internal (I) & do-calculus as the language of ablation claims; ablation and patching as formally distinct interventions; transportability for cross-model transfer \\
\midrule
Neuroscience & Internal (I) & Lesion versus stimulation, double dissociation as the test for specificity, and convergence across modalities as the unit-of-analysis criterion \\
\midrule
Genetics & Internal (I) & Epistasis for non-additive component interaction, rescue experiments for reversibility, and sensitivity analysis for unmeasured confounders \\
\midrule
Medical microbiology & Internal (I) & Graded presence in place of presence-absence, evidence gathered at the level the claim is pitched at, and partial satisfaction as sufficient for a causal conclusion \\
\midrule
Pharmacology & External (E) & Dose-response as a minimum bar for magnitude claims, affinity versus efficacy, and system reserve as the reason ablation understates a component's role \\
\midrule
Mechanistic interpretability & Interpretive (V) & Description versus explanation, faithfulness versus understanding, and component identity versus component role \\
\bottomrule
\end{tabular}
\end{table}


% ============================================================
\section{Case Studies}
\label{sec:case-studies}

We apply the framework to sixteen published mechanistic claims, drawn from fifteen papers, evaluating each through all five validity lenses. We audit \emph{claims}, not papers: a verdict characterizes the current evidential status of a claim, aggregating all published evidence to date, and is not an assessment of any single paper's quality. Because verdicts are graded against an idealized standard, a low tier marks distance from that ideal, not a flawed study. Quantitative per-paper scoring is the subject of a companion paper; here we keep the analysis qualitative and trace each verdict to specific missing or present evidence.

Every claim below was re-derived against all 36 criteria from the origin paper read in full, with the post-origin literature checked against primary sources. Two conventions follow from that. Where a criterion is Untested we distinguish \emph{unattempted} from \emph{not applicable}, because only the first is a gap the authors could have closed. And where the claim as the field reads it is stronger than the claim the origin paper made, the row says which one is being scored---this matters most for sparse-autoencoder features and for probing classifiers, whose origin papers hedge in the direction their critics later took.

\subsection{Validity audit}

\Cref{tab:case-studies} presents the verdict assignments. The tier system discriminates: the sixteen claims spread across five labels, and each difference is traceable to a named criterion rather than to an overall impression. Two claims from the same paper land at opposite ends---induction heads reach Triangulated for token copying and Disconfirmed for general in-context learning---which is the clearest demonstration that the unit of assessment is the claim and not the paper.

\clearpage
{\footnotesize
\renewcommand{\arraystretch}{1.1}
\begin{longtable}{@{}>{\raggedright\arraybackslash}p{2.2cm}>{\raggedright\arraybackslash}p{2.2cm}>{\raggedright\arraybackslash}p{3.4cm}>{\raggedright\arraybackslash}p{7.4cm}@{}}
\caption{Verdict assignments for sixteen mechanistic claims, sorted by evidential status (see \Cref{tab:verdicts} for tier definitions). Scope records the task and system within which the verdict holds; a verdict does not travel beyond it. Criterion identifiers (C1--V5) are defined in \Cref{sec:criteria}; section references carrying no citation are to the origin paper named in the row.}
\label{tab:case-studies}\\
\toprule
\textbf{Verdict} & \textbf{Circuit} & \textbf{Scope (task, system)} & \textbf{Evidence characterization} \\
\midrule
\endfirsthead
\multicolumn{4}{@{}l}{\emph{Table \thetable, continued from previous page}}\\
\toprule
\textbf{Verdict} & \textbf{Circuit} & \textbf{Scope (task, system)} & \textbf{Evidence characterization} \\
\midrule
\endhead
\endfoot
\bottomrule
\endlastfoot
Mechanistically Supported & Fourier mechanism for modular addition \citep{nanda2023progress} & Modular addition mod 113; a one-layer ReLU transformer trained from scratch, scope declared by the authors & The algorithm is given in closed form, and the origin's Eq.~2 fixes the logit-axis structure of the neuron--logit map with no free parameters. Coverage stops short of the model: $W_{\mathrm{in}}$ receives no closed form, the neuron-side directions are free, and 79 of 512 neurons are unaccounted for. The mechanism recurs across moduli, seeds and data fractions (E2). A rival algorithm exists for the same task in structurally similar networks \citep{zhong2023clock} and has never been run against this model (I5); no double dissociation has been attempted (I6) \\
\midrule
Triangulated & Induction heads: token copying \citep{olsson2022context} & In-context token copying; causal evidence at origin on 1--6 layer models at $d_{\mathrm{model}} = 768$, extended to 70B by later work & Causal ablation reaches 70B---removing the top 10\% of induction heads collapses 8-shot accuracy from 80.57\% to 2.50\%, against 79.13\% for matched random ablation \citep{yang2025unifying}---conditional on a directly copyable label appearing in the demonstrations. The construct is defined on repeated \emph{random} sequences, so the claim is distribution-free by construction (E2), and \citet{feucht2025dual} supply a double dissociation between token-level and concept-level induction \\
\midrule
Mechanistically Supported & Greater-Than \citep{hanna2023greater} & Year-span prediction (``The \textless noun\textgreater{} lasted from the year XXYY to the year XX\_\_''); GPT-2 small, one released checkpoint & Necessity by sign reversal: routing the counterfactual through the circuit takes probability difference from $+81.7\%$ to $-36.6\%$ (\S3.2), and the complement test recovers 89.5\% of the unpatched model. The comparison operation itself is unresolved---``We cannot provide a conclusive answer'' (\S4)---and the one structural hypothesis tested causally is rejected (\S4.1). Specificity is tested and adverse: the circuit fires where less-than is required (\S5). Cross-model recurrence was tested later and failed \citep{xu2026pattern} \\
\midrule
Mechanistically Supported & Copy Suppression \citep{mcdougall2023copy} & Next-token prediction on OpenWebText; attention head L10H7 in GPT-2 Small & The mechanism is read off the weights before any ablation---95.72\% of QK rows have the diagonal largest, 84.70\% of OV columns have it among the ten most negative---against a baseline built for GPT-2's tied embeddings (\S3.1--3.2). Coverage is measured over the pretraining distribution rather than a task (E2) and runs 45--95\% depending on the metric, with 76.9\% the headline KL figure. No second mechanism has been shown intact under a copy-suppression-removing ablation (I6) \\
\midrule
Mechanistically Supported & Superposition \citep{elhage2022superposition} & Sparse feature reconstruction and $\mathrm{abs}(x)$; four toy ReLU networks on synthetic sparse data, all trained by the authors & The mechanism is derived rather than inferred: closed-form losses for each candidate weight configuration, the transition confirmed analytically as first order, the pentagon radius predicted from the bias and observed. Planted features supply a known-positive control (M5) and the linear model a control that cannot exhibit the phenomenon (M2). The scope is declared in the title, and the three criteria that would test travel beyond it (E2--E4) rest on later work by other groups \\
\midrule
Mechanistically Supported & Steering Vector (refusal direction) \citep{arditi2024refusal} & Refusing harmful instructions in a chat template; 13 chat models across 5 families, 1.8B--72B & Necessity and sufficiency by direct intervention across all 13 models, with off-target cost measured on six benchmarks and reported where it is adverse (I4). The adversarial-suffix analysis is a confirmed novel prediction rather than a fragility: the suffix works by suppressing the direction (\S5.1). The cap is mode-match---the extraction contrast is a harmfulness contrast, and \citet{zhao2025harmfulness} later separate a harmfulness direction from the refusal direction (V2). No intervention-strength sweep appears anywhere in the paper (E5) \\
\midrule
Mechanistically Supported & Global Workspace / J-space \citep{gurnee2026verbalizable} & Five declared workspace properties, no single task; four closed-weight Claude models, sizes undisclosed & Necessity established four times, three of them with matched controls, and specificity tested three ways, including two within-stimulus dissociations that move explicit report while leaving continuation untouched (I4). The circularity is confined to verbal report, is flagged by the authors, and is answered by a control on concept vectors derived without the lens. The cap is that nearly every intervention edits along lens directions: one evidence family, one laboratory, one model family, and no open-weight replication (C3, E4) \\
\midrule
Causally Suggestive & Othello World Model \citep{li2023othello} & Predicting legal next moves in Othello transcripts; Othello-GPT, 8 layers, trained from random initialization; primary causal evidence on the synthetic checkpoint & Board state is recovered by nonlinear probes at 1.7\% error where linear probes fail at 20.4\% (\S3.1), and 2,000 interventions move predictions onto the counterfactual legal-move set (errors 0.12 and 0.06 against null baselines of 2.68 and 2.59), including on boards unreachable by legal play (\S4.2). \citet{nanda2023emergent} later recover a \emph{linear} representation by re-basing the probe target. The label is never contrastively tested at origin; \citet{vafa2024evaluating} build the contrast, and the synthetic checkpoint passes it where the championship checkpoint fails \\
\midrule
Mechanistically Supported & Successor Heads \citep{gould2023successor} & Successor prediction over eight ordinal token classes; behavioral recurrence across GPT-2, Pythia and Llama-2 (31M--12B); mechanistic case study on Pythia-1.4B L12H0 & Convergent evidence for the mod-10 features from three methods run to exclude a sparse-autoencoder artifact (\S3.3), and the behavioral decomposition holds across five ablation regimes (their Appendix K). Necessity is the gap the authors name: ``we did not explicitly demonstrate that successor heads are necessary for incrementation'' (\S4). Specificity fails openly---acronym 23.8\% and greater-than 18.9\% of winning cases. Mechanistic recurrence is single-model, and \citet{xu2026pattern} find no shared causal structure across models \\
\midrule
Causally Suggestive & Docstring Circuit \citep{heimersheim2023docstring} & Predicting the next argument name in a Python docstring; a 4-layer attention-only toy transformer, 50 prompts per experiment & The authors name the mechanism ``Docstring Induction,'' name the simpler line-number rival, suppress it by prompt design at a measured cost ($\sim$75\% to 56\% model accuracy), and conclude both are implemented---so what is open is apportionment, not identity (I5, V3). The circuit recovers 42\% against the full model's 56\%. Specificity is unattempted (I4), and a $10^6$-pair stress test finds the circuit diverging from the model on benign in-distribution inputs \citep{uitdebos2024adversarial} \\
\midrule
Causally Suggestive & Gender Bias Circuits \citep{vig2020gender} & Bias as a probability ratio between two candidate continuations; GPT-2 distil--xl plus five further model families & Interchange intervention on heads and neurons against a randomly initialized control that separates cleanly (M2), giving sparsity (10 of 144 heads match the full effect), synergy, and decomposition. Necessity is untestable as designed---every estimand is a substitution, never a removal (\S3.2.2). No capability measure is defined anywhere in the paper, so specificity has no second quantity to register on (I4), and whether bias separates from gender competence is presupposed rather than tested (C4) \\
\midrule
Proposed & Probing Classifiers \citep{belinkov2022probing} & No single task: a claim form over a model's intermediate representations; the anchor is a review with no experiments, so the evidence sits in the works it cites & Probing carries no intervention of its own, and the causal follow-ups disagree---removal degrades the original task in one study and not in another \citep{elazar2021amnesic}. Probe accuracy does not discriminate between information the model uses and information merely present: \citet{hewitt2019designing} report 92.8\% on a control task with randomized per-word-type labels, and control datasets recover properties that are non-discriminative for the original task. Decodability is a legitimate result at this tier; the strong reading, that accuracy shows use, is what fails \\
\midrule
Causally Suggestive & IOI Circuit \citep{wang2023interpretability} & Indirect object identification, 15 templates in both name orders; GPT-2 small, 26 (head, position) pairs in seven classes & Scored at the most lenient reading, on which these heads carry the computation. Rival sets bear on uniqueness, which that reading leaves unasserted: the origin's own na\"ive circuit reaches comparable faithfulness (\S4.3), and \citet{chen2026circuits} exhibit two IOI subgraphs at 100\% accuracy with 4.1\% edge overlap. E1 caps the claim. The 87\% figure is an average under one protocol: \citet{miller2024transformer} move the same quantity from below 0\% to well over 100\% across six methodological choices, faithfulness splits by template order (M6), and \citet{uitdebos2024adversarial} find the circuit diverging from the model on benign in-distribution inputs (E2) \\
\midrule
Disconfirmed & Knowledge Neurons \citep{dai2022knowledge} & Cloze recall of a relational fact with a single-word tail; BERT-base-cased only, roughly four FFN neurons per fact & Suppression and amplification move correct-answer probability by $-29.03\%$ and $+31.17\%$ against a count-matched baseline's $-1.47\%$ and $-1.27\%$, which the authors present as a proof of concept. Off-target effects are measured at origin and point both ways: erasure spares other relations, while fact updating damages unrelated relations \emph{more} than editing random neurons does (their Table 6, 7.2 against 4.3). \citet{niu2024knowledge} then move non-factual linguistic phenomena with the same machinery---the construct does not separate from its neighbours (C4), and the storage claim outruns the evidence (V2) \\
\midrule
Disconfirmed & SAE features as canonical units \citep{cunningham2023sparse, leask2025canonical} & Method-level claim, no single task; Pythia-70M and 410M at origin, GPT-2 small and Gemma 2 2B for the canonicity tests & Scored as the field's strong reading; the origin paper hedges in its critics' direction. Large-dictionary latents are compositions of \emph{coarser} components matching smaller-dictionary latents, with the sparsity penalty named as the cause \citep{leask2025canonical}. The origin's headline interpretability metric returns similar values on randomly initialized transformers \citep{heap2025sparse}; cross-seed agreement falls from 94\% at 3{,}072 latents to 30\% at 131K \citep{paulo2025sparse}; and features localize causal variables no better than raw neurons \citep{mueller2025mib} \\
\midrule
Disconfirmed & Induction heads: general in-context learning \citep{olsson2022context} & General in-context learning; transformers of any size (asserted in the abstract), with the causal evidence confined to models under 42M parameters & At 70--72B the heads carrying abstract in-context reasoning are disjoint from induction heads ($r=0.11$ against prefix-matching, $r=0.86$ against function-vector scores) \citep{webb2025emergent}; suppressing induction-head formation preserves abstract in-context learning on 13 of 21 tasks \citep{sahin2025incontext}; with function-vector heads preserved, induction ablation above 1B matches random ablation \citep{yin2025which}. What fails is few-shot task accuracy, the metric the field substituted---on the origin's own token-loss-difference measure induction heads remain dominant above 345M \citep{yin2025which}. The assumption under attack is one \citet{olsson2022context} identify themselves: their Argument~6 concedes the case ``is stronger for small models than for large ones'' and bridges the gap by continuity, which is the premise the scale evidence denies \\
\end{longtable}
}


\subsection{Worked examples}

We describe seven case studies in detail to illustrate how the framework produces its verdicts. These assessments hold claims against an idealized standard; a low verdict marks distance from that ideal, not a flawed study.

\paragraph{Induction Heads (Triangulated and Disconfirmed).} \citet{olsson2022context} identify a two-head composition mechanism for in-context copying: previous-token heads attend to the token before a repeated sequence, and induction heads use this signal to predict the next token. The paper advances two claims of different reach, and the evidence has since separated them.

The narrower claim---induction heads implement in-context token copying---reaches Triangulated. Causal ablation now extends to 70B: removing the top 10\% of induction heads in Llama2-70B collapses 8-shot accuracy from 80.57\% to 2.50\%, against 79.13\% for a matched random-head ablation \citep{yang2025unifying}, which satisfies E4 with an intervention rather than a correlation. That number carries a condition worth stating, because it bounds the mechanism rather than the claim: when the query's label is absent from the demonstrations the ordering reverses and removing previous-token heads costs more than removing induction heads. The origin's own causal evidence stops at six-layer models, so the scale here is later evidence, not the origin's. The construct is defined on repeated sequences of uniformly random tokens, which makes the claim distribution-free by construction (E2), and \citet{feucht2025dual} later split it cleanly: mean-ablating concept induction heads collapses word-level translation while leaving nonsense copying intact, and the converse holds---a double dissociation obtained three years after the claim (I6).

The broader claim---that induction heads are the mechanistic source of \emph{general} in-context learning---is Disconfirmed. At 70--72B the heads carrying abstract in-context reasoning form a disjoint set from induction heads, correlating $r=0.11$ with prefix-matching scores and $r=0.86$ with function-vector scores \citep{webb2025emergent}. Suppressing induction-head formation during training leaves abstract in-context learning intact on 13 of 21 tasks \citep{sahin2025incontext}. And once function-vector heads are preserved, ablating induction heads above 1B is comparable to ablating random heads \citep{yin2025which}, though the same models continue to show induction dominance on the token-level metric \citet{olsson2022context} originally used---so the two claims diverge by outcome measure as well as by scale.

This is a scope refutation the authors anticipated. Their Argument 6 states that ``above some size, non-induction composition heads could together account for more of the phase change and in-context learning improvement than induction heads do.'' No replication was attempted and failed; the narrow mechanism replicates, and the reach asserted for it does not. One paper, two claims, two tiers---the same structure as the workspace case, and the reason the audit scores claims rather than papers.

\paragraph{IOI Circuit (Causally Suggestive).} \citet{wang2023interpretability} identify 26 attention heads grouped into seven classes for indirect object identification. Sufficiency is the headline: mean-ablating everything outside the circuit leaves a logit-difference gap of 0.46, ``or only 13\% of $F(M) = 3.56$,'' which is the 87\% figure. Necessity at the component level is weaker than the number suggests, and the paper says so---knocking out all three Name Mover Heads costs ``only 5\% drop in logit difference'' (\S3.4), and the minimality scores are conditional on a co-ablated set chosen to maximize them.

Three findings move the claim off Mechanistically Supported, and none of them is a shortage of evidence. First, rival specifications explain the behavior as well: the paper's own na\"ive baseline circuit reaches comparable faithfulness (\S4.3), and \citet{chen2026circuits} exhibit two IOI subgraphs at 100\% task accuracy sharing 4.1\% of their edges. Second, the faithfulness number is a property of the protocol as much as of the circuit---\citet{miller2024transformer} vary six methodological choices and the same quantity ranges from below 0\% to well over 100\%, with the variant that ``best represents the hypothesis of Wang et al.'' scoring a median well over 100\%. Third, the circuit diverges from the model on inputs drawn from the origin task itself: over $10^6$ benign clean/corrupted pairs the model-circuit KL has mean 5.15 and maximum 14.64 \citep{uitdebos2024adversarial}. Specificity is tested rather than untested, with results that depend on the quantity measured---78\% head overlap with the Colored Objects circuit in GPT-2 Medium \citep{merullo2024circuit} against near-zero cross-task faithfulness in GPT-2 small \citep{hanna2024faith}.

Each claim is scored at its most lenient reading, and the lenient reading here is that these 26 heads carry the indirect-object computation. Rival subgraphs are compatible with it, since a second sufficient set establishes that the mechanism is not unique and uniqueness belongs to the stronger reading. Method disagreement (E1) caps the claim instead, and no reading escapes that. The localization is real and reproducible, and which subgraph constitutes \emph{the} mechanism is open.

\paragraph{Probing Classifiers (Proposed).} Probing \citep{belinkov2022probing} trains a classifier on activations to test whether a property is decodable. The core limitation is one the survey itself names as its headline shortcoming: the probing framework ``may indicate correlations between representations \ldots\ and linguistic property $z$, but it does not tell us whether this property is involved in predictions of the model'' (\S4.3). Probe success establishes decodability, and high accuracy is equally consistent with genuine use, incidental separability in a high-dimensional space, and encoding of a confound. Three designed tests sharpen the ceiling. Control tasks show that popular probes are not selective---an MLP probe reaches 97.3\% on part-of-speech against 92.8\% on a task with randomized per-word-type labels, so most of the accuracy is the probe's own capacity \citep{hewitt2019designing}. Control datasets show that properties are recovered above chance even when they are non-discriminative for the original task. And amnesic probing shows that ``conventional probing performance is not correlated to task importance'' \citep{elazar2021amnesic}.

The ceiling is therefore sharper than ``no internal validity,'' and the survey records why: an earlier intervention study on subject--verb agreement concluded that probes \emph{do} identify features the model uses, and \citet{elazar2021amnesic} concluded the opposite, with both results standing. What probe accuracy cannot do is tell you which case you are in. Probing \emph{with} causal follow-up can reach Causally Suggestive; with control tasks and cross-method convergence it can reach Mechanistically Supported. Probing is not dismissed---its evidential reach is bounded, and the bound has been measured.

\paragraph{Gender Bias Circuits (Causally Suggestive).} \citet{vig2020gender} apply causal mediation analysis to gender bias and report three properties of the effect: it is sparse (ten of 144 attention heads reproduce the effect of intervening on all of them), synergistic (for neurons, the concurrent indirect effect far exceeds the sum of independent ones), and decomposable into direct and indirect parts. The causal evidence is real and controlled: a randomly initialized GPT-2 small is carried through the analysis as a matched negative control and separates at both levels, which is why the sparsity result is attributed to training rather than to architecture.

Two structural features of the design bound how far it reaches, and both are visible in the methods rather than in the results. The estimands are substitutions: the natural direct effect holds the mediator at its original value and the natural indirect effect sets it to its counterfactual value (\S3.2.2). Nothing is removed, so necessity cannot be expressed in this design---a criterion untestable in the setting as run, not a gap the authors declined to fill. And the outcome is a probability ratio between two candidate continuations, with no accuracy, perplexity, or capability measure defined anywhere in the paper, so there is no second quantity on which an off-target effect could register (I4).

The interesting failure is upstream of both. The name ``gender bias circuit'' presupposes that bias separates from gender competence---that one could intervene on the biased prediction without disturbing the model's ability to resolve a gendered pronoun that is licensed lexically rather than stereotypically. The paper never localizes a gender-competence mechanism, so no overlap statistic exists; the definitionally gendered professions that would carry such a test are excluded from the total-effect calculation by design. Discriminant validity here is untested (C4), which is a different verdict from failed, and the difference is exactly the one the framework exists to enforce.

The literature has since begun testing it, and disagrees. \citet{veloso2026gknow} run the two-outcome design on Llama-3.1-8B and Olmo-7B and report that ``gender bias and factual gender are severely entangled on the level of both circuits and neurons'': ablating the top-50 stereotype-attributed neurons in Llama costs factual-gender accuracy 91.49 to 79.86, while ablating only the neurons \emph{unique} to the stereotypical set moves nothing at all---the sharp form of a missing therapeutic window, since the non-shared components are inert. Against that, \citet{yu2025gender} report reduced bias with capabilities preserved on five models. Neither study uses GPT-2, neither uses mediation analysis, and both operate on lexical rather than grammatical gender, which English does not mark. The honest state is contested rather than settled, and the claim sits where its own evidence puts it: causal intervention with a matched control, no necessity primitive, no specificity measure, and a construct whose central presupposition is still open.

\paragraph{SAE features as canonical units (Disconfirmed).} A sparse-autoencoder feature labeled as a concept---``this feature is deception,'' ``this feature is the Golden Gate Bridge''---claims that the model represents that concept along a recovered direction, and the strong version of the claim adds that such directions are the model's units of analysis. The strong version is what fails, and it is worth separating from what \citet{cunningham2023sparse} assert. Their claims are comparative throughout---features ``more interpretable and monosemantic than directions identified by alternative approaches''---and they anticipate non-canonicity in their own results, noting that the absence of a knee in the sparsity--accuracy tradeoff ``provides some evidence that there is not a single correct way to decompose activation spaces into a sparse basis.''

Three designed tests then settle the strong reading against it. \citet{leask2025canonical} show that large-dictionary latents are compositions rather than atoms, and the direction of decomposition is the informative part: an ``Einstein'' latent resolves into \emph{coarser} meta-latents---``scientist,'' ``Germany,'' ``famous person''---which match latents in \emph{smaller} dictionaries. That is not feature splitting, which runs finer with scale; the authors name the cause as the sparsity penalty, which ``incentivizes larger SAEs to learn compositions of latents rather than atomic latents.'' It is a discriminant-validity (C4) failure, because the labeled unit is a composite of its neighbors. \citet{heap2025sparse} then remove the ground under the origin's headline metric: automated interpretability scores on SAEs trained over \emph{randomly initialized} transformers are similar to those from trained models, so the instrument does not distinguish model structure from dictionary structure. And feature identity depends on scale: \citet{leask2025canonical} report 94\% cross-seed agreement at 3{,}072 latents on GPT-2 (their Figure 16), while \citet{paulo2025sparse} find only 30\% shared features at 131K latents on Llama 3 8B. The two are measured on different models with different metrics, so the pair bounds the problem rather than tracing a single curve.

The causal comparison points the same way, though it rebuts a field expectation rather than an origin claim: \citet{mueller2025mib} find SAE features localize causal variables no better than raw neurons, a comparison \citet{cunningham2023sparse} never ran---their causal section compares against principal components and a non-sparse dictionary only. A feature-clamping intervention can still lift a \emph{specific} feature to Causally Suggestive. What is Disconfirmed is the generic claim, and the phrase ``features may be dictionary properties rather than model properties'' is now a measured result rather than a worry.

\paragraph{Steering Vector (Mechanistically Supported).} A steering vector claims that adding a direction to the residual stream drives a target behavior. Unlike a feature or a probe, this is interventional by construction---a $\mathrm{do}$-operation on the activations---so it clears the causal bar the observational cases cannot. The refusal direction \citep{arditi2024refusal} is the strongest instance: ablating it removes refusal, adding it induces refusal, and the effect replicates across 13 chat models in five families spanning a 40$\times$ parameter range. Off-target cost is measured rather than assumed, on six benchmarks plus cross-entropy loss on three corpora, and it is reported where it is adverse---TruthfulQA falls for all thirteen models.

One criterion the method class advertises is absent from this instance. Steering admits a dose-response by scaling the vector, but no such sweep appears here: activation addition is defined at unit magnitude and directional ablation is a projection with no magnitude parameter at all, and no appendix varies either (E5). That is a cheap experiment left unrun, not a limitation of the approach.

The claim caps on interpretive rather than internal validity. A direction that is a sufficient behavioral \emph{lever} is routinely narrated as \emph{the representation} the model uses for the concept---a description-mode mismatch (V2), the affinity-versus-efficacy confusion pharmacology names. The origin supplies the evidence for its own cap: the direction is extracted as a difference in means between harmful and harmless \emph{instructions}, a contrast over input content, and is then labeled a refusal direction without a test of which construct it carries. \citet{zhao2025harmfulness} run that test and find a harmfulness direction distinct from the refusal direction, with different steering consequences---steering along harmfulness makes the model read harmless instructions as harmful, while steering along refusal elicits refusal without moving the judgment. The lever and the representation come apart, which is what a lever-not-mechanism account predicts. The paper is itself more careful than its reception: it calls the work ``more of an existence proof that such a direction exists'' and writes ``refusal direction'' in scare quotes throughout. The live boundary condition is the word \emph{single}---independent extraction methods recover multiple directions, and a rank-4 projection has since been reported to give a better utility--safety balance than the rank-1 one.

\paragraph{Global Workspace / J-space (Mechanistically Supported).} Anthropic's Jacobian lens assigns every vocabulary token a direction in the residual stream, derived from the expected Jacobian of the final-layer residual with respect to a layer's activations; the ``J-space'' is then the set of points expressible as a sparse nonnegative combination of those directions \citep{gurnee2026verbalizable}. The paper narrates it as a ``global workspace'' holding the concepts the model reasons over. The mechanistic core is unusually strong. Necessity is established four times with matched controls---J-space ablation drives multi-hop accuracy to near zero, broadcast-head ablation drops injected-concept reporting from 0.54 to 0.09 against a layer-matched random-head control, and two further token-level ablations move the behaviors they target.

Specificity is tested, and tested well: a fourteen-task battery separates shallow classification and factual recall, which survive heavy ablation, from inference-grounded generation, which collapses. Two within-stimulus designs are stronger than the battery. A single Spanish-to-French swap flips explicit report and flexible inference while leaving continuation and anomaly detection untouched---with the word \emph{Spanish} present in the lens readout in all four conditions, so the dissociation is causal rather than representational. A count swap moves a reported character count from 46 to 65 while leaving the line-wrap point unchanged.

The claim caps at Mechanistically Supported for a reason the framework makes precise, and it is not specificity. Nearly every intervention edits along lens directions, so the evidence is one family (C3); one of three lens-independent checks, mean-difference activation patching of arithmetic intermediates, agrees on the causal depths. And every result comes from four closed-weight models in one family at one laboratory, with no sizes, layer counts, or architectural parameters disclosed and no open-weight replication (E4). Circularity is a narrower charge than it first appears: the lens is derived from causal effects on output tokens, so a relationship to verbal report is expected---the authors say so twice---and they build the control, deriving concept vectors without the lens and finding that the J-space component carries 6--7\% of their variance yet drives 59\% of report swaps against 5\% for the 93\% outside it. The paper's deception result has its contrastive foil in the model organisms rather than in the widely quoted blackmail transcript, which is a single qualitative readout with no baseline; the authors decline the monitoring claim outright.

\subsection{Recurring validity gaps}

Several patterns emerge from evaluating these claims side by side.

\paragraph{Unattempted criteria (M7, I6, I10).} Necessity and sufficiency track together across the set rather than forming a gap: where I1 is established, I2 usually is too, and where I1 is partial, so is I2. The systematic absences are elsewhere, and they are near-total. Selection correction (M7) is Untested in all sixteen claims, even though every one of them states its candidate pool---144 heads, 512 neurons, 56 frequencies, $|I| \times L$ candidate directions---and then reports the selected item's score as its property. Copy suppression shows the shape of the cost most clearly, in a paper that discloses every step. L10H7 is chosen as the most negative head in GPT-2 Small; the 80\% behavioral figure is computed on completions sampled from the top 5\% by L10H7's own effect; the QK ablation retains the top 5\% of predicted source tokens; and both faithfulness appendices evaluate on the top 5\% again. Each choice is stated with its reason, which is more than most papers offer. None is recomputed on an unselected slice, so nothing bounds how much of either headline number belongs to the head and how much to the selection. What the criterion asks for is the held-out estimate rather than a multiplicity correction, and it gates no tier: whether a candidate pool was disclosed bears on the size of an effect, not on whether the mechanism is there. Double dissociation (I6) and rescue reversibility (I10) are Untested in fifteen and thirteen of sixteen respectively, and in most of those the missing experiment is cheap---ablation procedures that cache clean activations make corrupt-then-restore one extra forward pass. The instruments exist and sit unused. Confounding sensitivity has a developed methodology in causal inference---E-values \citep{vanderweele2017sensitivity}, omitted-variable bounds \citep{cinelli2020making}, and the $\Gamma$ bounds of \citet{rosenbaum2002observational}---drawn from the same literature the field already borrows its interventional vocabulary from. What these four criteria share is that nobody has tried them.

\paragraph{Attempted but unconfirmed criteria (M4, I5, C4, M3, I3).} A second group is attempted almost everywhere and confirmed almost nowhere. Calibration (M4) is attempted in all sixteen claims and Confirmed in none; fourteen land at Partially confirmed, because a metric is reported in interpretable units and never checked against its nearest false positive. Rival mechanism exclusion (I5) is attempted in fifteen and Confirmed in none. Discriminant validity (C4), stability (M3) and minimality (I3) repeat the pattern at eleven to fourteen attempts each. Minimality is the clearest case, because the missing work is procedural rather than conceptual: circuit discovery has a complexity-theoretic treatment \citep{adolfi2025complexity} and algorithms with provable guarantees over a hierarchy running from quasi-minimal through locally- and subset-minimal to cardinally-minimal \citep{hadad2026formal}, and no audited claim establishes minimality at any level of it. This group is the more tractable of the two, because the experiments are already being run and stop short of the standard.

\paragraph{Method-conditional results.} The IOI circuit's headline number is a property of the protocol as much as of the circuit: varying six methodological choices moves the same faithfulness quantity from below 0\% to well over 100\% \citep{miller2024transformer}, and the direction of the change is not even consistent---resample ablation scores systematically lower than mean ablation at the node level, while the edge-level variant closest to the original hypothesis scores highest of all. Ablation type is part of the causal claim, not an implementation detail. The framework captures this through E1 (intervention reach), which requires agreement across ablation methods, and the pattern recurs: successor heads survive a five-regime recomputation on the headline category while their minor categories move sharply, and the greater-than circuit's relative faithfulness flips sign under one ablation family.

\paragraph{Toy-model ceiling.} Two claims in the set are validated inside a scope their authors declared in the title, and both illustrate how the framework handles that. The Fourier mechanism for modular addition reaches Mechanistically Supported on three genuinely distinct families---periodicity in activations, Fourier structure in the weights, and frequency ablation---plus progress measures that track the mechanism forming before the behavior appears. Superposition reaches Mechanistically Supported on evidence that is derived rather than inferred: closed-form losses for each candidate weight configuration, an analytically located first-order phase transition, and planted ground-truth features that make a known-positive control available at every point. What neither has is travel. Superposition's E2, E3 and E4 all rest on later work by other groups, and its input distribution is one synthetic family throughout. A claim validated within a declared scope is a legitimate result, not a failure (V5); the gap between toy-model proof-of-concept and real-model confirmation remains the field's central challenge.

\paragraph{No circuit reaches Validated.} The Validated tier requires all five validity types to be addressed: construct, measurement, internal, external, and interpretive. No published circuit currently meets this bar. Induction heads come closest on the narrow claim---construct, internal and external validity are all strong, with a double dissociation now on the record---but reliability is reported nowhere in the origin paper despite the authors having trained all 35 models themselves, and every head-evaluator score is a single point estimate. The Validated tier is a concrete checklist, and the gap to it is enumerable. The gap between Triangulated and Validated is primarily measurement-validity infrastructure (bootstrap stability, seed variance, baseline separation for the metrics used to evaluate the circuit) and interpretive discipline (explicit level declaration, checks on unlicensed labels). These are tractable additions to existing experimental practice, not philosophical impossibilities.

To demonstrate reachability, consider what it would take for the Fourier mechanism for modular addition---currently Mechanistically Supported, since no double dissociation has been attempted (I6)---to reach Validated. It supplies a closed-form target mechanism with measured residuals rather than an exact ground truth: the algorithm is given in closed form and the neuron--logit map's logit-axis structure follows from it with no free parameters, while 79 of 512 neurons fail the polynomial fit, two of four attention heads are assigned a role the authors introduce with ``we speculate,'' and the headline logit reconstruction is 95\% rather than exact. The remaining gaps are: (1) M1--M7: report bootstrap or interval estimates on the fraction-of-variance-explained figures that carry the mechanism claim, which are currently single point estimates, and perturb the analysis-side thresholds---what counts as a nontrivial Fourier coefficient, the rank-10 truncation, the 85\% cutoff that selects 433 neurons---none of which is varied; (2) V1--V5: declare the description mode, and audit whether ``fully reverse engineer'' is licensed by evidence that leaves 15\% of neurons undescribed. Each is a straightforward addition to the existing analysis---no new experiments required, only systematic reporting of calibration checks and interpretive discipline. Alternatively, a compiled transformer \citep{lindner2024interp} has synthetic ground truth by construction: the circuit \emph{is} the mechanism, because the weights were generated from a known program. Running the full M1--M7 suite and V1--V5 audit on such a model would produce the first Validated claim, proving the tier is achievable and calibrating the framework against a known answer.

\paragraph{Interpretive inflation.} Labels like ``world model,'' ``knowledge neuron,'' and ``deception feature'' carry theoretical implications beyond what the evidence supports. The diagnostic is that the label is never contrastively tested: \citet{li2023othello} define ``world model'' explicitly and then run no experiment separating it from a decodable state summary that is causally used for move prediction, and it took a later paper building that instrument to settle the question in the authors' favor for one of their two checkpoints. The framework identifies where labels exceed evidence through V4 (unlicensed labeling) and V5 (scope declaration), and the remedy is a designed contrast rather than a milder word. A name is licensed when it redescribes the measurement, and unlicensed when it picks out a construct that a nearby alternative would satisfy just as well with nothing run to separate them. Anthropomorphism is the salient case rather than the criterion: ``Name Mover'' carries no intentional term and still asserts a function, while a label the authors mark as speculation asserts nothing they have not flagged.

\paragraph{Roles named without being filled.} Nine of the sixteen claims name at least one component by the role it plays while stating that no account exists at the level below. The Fourier mechanism assigns two of four heads an amplification role that the authors flag as speculation. The greater-than circuit's comparison step is one its authors say they ``cannot provide a conclusive answer'' about. IOI leaves four of seven head classes without a weight-level account. By Craver's test these are the innocuous kind of gap: in every one of the nine the origin paper states the gap itself, so the provisional status is on the record. What the completeness dimension tracks is whether it stays there. A role name reused as settled, without the gap restated, has moved from placeholder to filler with no new evidence intervening---and that transition happens in the citing literature rather than in the origin paper.

\paragraph{Untested discriminant validity.} Three claims in the set name a construct whose separation from a neighbor is presupposed and never measured. ``Gender bias circuit'' presupposes that bias separates from gender competence, and \citet{vig2020gender} localize only the first, excluding from their measurement the definitionally gendered items that would carry the test. ``Knowledge neuron'' presupposes that factual knowledge separates from other linguistic regularity, and \citet{niu2024knowledge} later show the same editing machinery moves both. ``SAE feature'' presupposes that a recovered direction is a unit of the model rather than of the dictionary, and \citet{heap2025sparse} show the instrument that established it cannot tell a trained transformer from a randomly initialized one. C4 is the criterion most often skipped and the one whose failure is least repairable by more of the same evidence: where the construct is not separable from its neighbor, the response is reformulation, not another experiment.

% ============================================================
\section{Discussion}
\label{sec:discussion}

The case studies reveal that the dominant validity gaps in published MI research are systematic, not idiosyncratic, and that they are not the gaps the field's own commentary names most often. Specificity is tested more than its reputation suggests, and where it has been tested the answer often disagrees across models or quantities. What is absent is uniform: no claim in the set corrects for the selections it makes, and none bounds the effect of an unmeasured confounder. Faithfulness results move with the intervention protocol, sometimes in both directions. Labels are asserted without the contrastive test that would separate them from a weaker description. These are structural features of a field that evaluates methods and artifacts without a common protocol for evaluating the claims built on them.

A principle implicit in the tier definitions deserves explicit statement: evidence accumulates \emph{within} a validity type, not across types.  No quantity of behavioral evidence can substitute for missing causal evidence, and no number of single-method causal experiments can substitute for cross-method convergence.  The dependency chain (construct $\to$ measurement $\to$ internal $\to$ external $\to$ interpretive) imposes a ceiling: a claim's tier is bounded by its weakest validity dimension, not raised by its strongest.

The framework addresses this by providing a structured validity profile for each claim---making gaps visible and the path to strengthening them concrete. An evidence card format operationalizes this as an atomic reporting unit: each card records which measurement bears on which criterion, whether the relevant calibration checks were performed, and what verdict the evidence supports. We propose that evidence cards serve as a standard reporting format for MI claims, analogous to CONSORT checklists for clinical trials. Organizations deploying MI-based monitors (deception probes, steering vectors, circuit-level classifiers) can evaluate the underlying features through the same criteria: Is the construct well-defined? Is the measurement reliable? Is the evidence causal? Does the finding generalize? A feature at Proposed should not be deployed as a monitor; a feature at Mechanistically Supported or above has the minimum evidential basis for cautious deployment with ongoing validation.

Two safety-relevant applications illustrate how the framework extends beyond retrospective auditing of published circuits.

\paragraph{Chain-of-thought faithfulness.}
The claim ``chain-of-thought reasoning is faithful'' is at least four claims, each requiring different evidence.  Does CoT predict the model's output (\emph{behavioral faithfulness}, testable with perturbation studies)?  Does it converge with other interpretability methods on the same internal structure (\emph{convergent faithfulness}, requiring multi-method evidence)?  Does it describe the functional roles actually implemented by internal components (\emph{role faithfulness}, requiring independent functional tests)?  Or does it map onto the model's computational graph (\emph{structural faithfulness}, requiring causal and structural evidence)?  \citet{lanham2023measuring} and \citet{turpin2024unfaithful} demonstrate that behavioral faithfulness often fails: models produce plausible reasoning chains that do not reflect their actual decision process.  Under Mechanistic Validity, ``CoT is faithful'' at the behavioral level is a Proposed claim (behavioral evidence only, no causal intervention); at the structural level, it would require at minimum Mechanistically Supported status across internal and construct validity criteria.  The framework makes the gap between these levels explicit rather than leaving ``faithful'' as an ambiguous umbrella term.

\paragraph{Deceptive alignment.}
Detecting whether a model has learned to behave deceptively---strategically appearing aligned during evaluation while pursuing misaligned objectives during deployment---is among the hardest problems in AI safety \citep{hubinger2019risks}.  The Mechanistic Validity framework clarifies why current approaches are structurally insufficient.  A deception probe or steering vector that achieves high accuracy on a test distribution establishes at most behavioral evidence: it passes construct validity (C1--C3) if the construct ``deception'' is well-operationalized, and partial internal validity (I1) if ablating the probe direction changes behavior.  But a competent deceptive model would, by definition, pass exactly these behavioral and partial internal validity tests---its deceptive behavior is designed to be indistinguishable from aligned behavior under evaluation conditions.  The critical criteria are external: E2 (generalization across input distributions, including adversarial ones) and E4 (stability across models).  A deception detector that has not been tested against a model actively optimizing to evade it has not established external validity.  Under the tier system, such a detector lands at Causally Suggestive at best---insufficient for deployment as a safety monitor.  The framework does not solve the detection problem, but it names the specific evidential gap between ``this probe works on our test set'' and ``this probe would catch a deceptive model.''

\paragraph{Safety evidence gaps at scale.}
More broadly, the dominant pattern in safety-relevant MI findings is instrumental evidence supporting structural conclusions.  Refusal direction ablation \citep{arditi2024refusal}, representation engineering for truthfulness, and feature-level safety classifiers all establish that a direction or feature is a sufficient behavioral lever---not that it is the mechanism.  The gap is systematic: safety monitoring requires at minimum object-level evidence (necessity plus sufficiency, I1--I2) and ideally subspace-level evidence (stability across distributions, E1--E6).  Most published safety-relevant findings, when audited against Mechanistic Validity criteria, land at Proposed or Causally Suggestive.  The framework provides a concrete upgrade path: which criteria are unmet, what evidence would meet them, and what tier the claim would reach with that evidence.

\paragraph{Minimum deployment standards.}
The safety applications above share a concrete implication: before a safety-relevant interpretability result is used in a deployment decision, what tier must it reach? We propose that the minimum standard is \emph{Mechanistically Supported} with the following criteria confirmed: I1 (necessity), I2 (sufficiency), I4 (specificity), E1 (intervention reach across at least two methods), and M2 (baseline separation against random and untrained controls). A result at Proposed or Causally Suggestive has not demonstrated that the measurement is trustworthy or that the effect is specific---deploying it as a monitor is the MI equivalent of prescribing a drug that passed Phase~I but not Phase~III. For adversarial settings (deception detection, jailbreak monitoring), the bar should be higher: E2 (generalization across input distributions, including adversarially chosen ones) is necessary because the threat model includes an agent optimizing to evade the monitor. These standards are not arbitrary---they follow from the dependency structure of the validity types: you cannot trust an external validity claim (``this generalizes'') without first establishing internal validity (``this is causal''), and you cannot trust a causal claim without first establishing measurement validity (``this metric works'').

\paragraph{The emergent abilities case and M4.}
The \citet{schaeffer2024emergent} emergent abilities mirage is the cleanest illustration of why M4 (calibration) belongs in the framework's measurement layer. The ``emergence'' of sudden capability jumps in LLMs was not a property of the models but of the evaluation metric: exact-match accuracy creates a sharp threshold where continuous improvement appears as a phase transition. Switching to a linear metric (token edit distance) dissolved the phenomenon entirely. M4 would have caught this: it requires that reported numbers correspond to meaningful quantities, and that the metric's response function be monotone and approximately linear in the range of interest. Had M4 been applied before the ``emergence'' literature took off, the question would have been ``does this metric create the pattern?'' rather than ``what explains this striking pattern?''---and the answer would have been immediate. The same diagnostic applies to any MI metric that produces sharp thresholds: if IIA jumps from 0.1 to 0.95 when a hyperparameter is tuned past a threshold, the first question should be whether the metric's response function is creating the jump, not whether the model's computation has a phase transition at that point.

\paragraph{Standards the field is ready for.}
The prescription here is not a matter of taste; it follows from the dependency structure and from precedent. Because external validity presupposes internal validity, which presupposes measurement validity, a field cannot trust a generalization claim it has not first shown to be causal and reliably measured---the ordering of evidence is forced, not chosen. Every mature science that makes mechanistic claims eventually codified this. Clinical medicine adopted CONSORT reporting and pre-registered trial protocols; psychology adopted registered reports; genetics adopted genome-wide significance thresholds and mandatory independent replication. Each standard began as a response to the same failure this paper documents---strong individual results outrunning the evidence that would license them. Mechanistic interpretability is now where those fields were before they codified: it has shared metrics but not shared validity criteria. We propose three practices as a matter of course: declare the description mode, report the evidence tier a claim reaches, and hold the strongest tiers to severe tests, whether preregistered or independently corroborated. These make the field's own implicit standards explicit rather than imposing an external one.

\paragraph{From framework to practice.}
Validity frameworks gain traction when they are socially enforced, not merely logically correct. We propose three adoption mechanisms, each modeled on a precedent from another field. First, \emph{validity statements}: authors declare which Mechanistic Validity tier their claims reach and why, analogous to the limitations sections that journals already require. Second, \emph{reviewer checklists}: conferences adopt a short Mechanistic Validity checklist (What description mode is declared? What tier does the evidence support? Is the claimed tier consistent with the evidence tier?), modeled on the CONSORT checklist in clinical trials. Third, \emph{evidence cards}: a standardized reporting format recording which measurement bears on which criterion, whether calibration checks were performed, and what verdict the evidence supports. These are tractable additions to existing practice, not new burdens. Why would individual researchers adopt them voluntarily? Because the framework converts an implicit weakness (``we didn't test X'') into an explicit strength (``our claim reaches Mechanistically Supported; here is the concrete evidence that would advance it to Triangulated''). Papers that use this framing give reviewers less to argue about and give follow-up studies a clear target.

Evidence on reporting checklists tempers the reviewer-checklist mechanism specifically. The IICARus trial randomized 1{,}689 manuscripts to an intervention designed to improve compliance with the ARRIVE guidelines, and no manuscript in either arm achieved full compliance \citep{hair2019randomised}. A checklist that authors complete without consequence changes what is claimed to be reported more than what is reported. Validity statements and evidence cards are proposed here as reporting units rather than compliance instruments, and their adoption depends on authors having a reason to prefer them: a card recording a claim at Mechanistically Supported, together with the specific evidence that would advance it, states a position an author would otherwise have to defend implicitly.

Several directions for future work follow naturally. Extending the case studies to larger models may reveal that criteria such as E4 (cross-model generalization) are harder to satisfy than the GPT-2 Small evaluations suggest. Automating the verdict assignment process---currently requiring human judgment in aggregating criterion-level results---is future work. Auditing a claim across training checkpoints would test whether a mechanism appears when the behavior does, evidence none of the sixteen currently carries. The models most circuits are found in release no intermediate checkpoints, so the question is unanswerable where the claim is made, and answering it on a suite that does release them, such as Pythia \citep{biderman2023pythia}, produces a result about a different model. Which artifacts a field releases therefore determines which criteria its claims can ever satisfy. And developing standardized prompt suites for specificity testing (I4) and double dissociation (I6) would address the two most common gaps identified across the sixteen case studies.

\section{Limitations}
\label{sec:limitations}

This work is a theoretical contribution. The case study verdicts are based on published evidence, not new experiments; running the framework's full metric suite with actual model interventions is future work. The current case studies evaluate the most widely-studied published circuits, which predominantly target GPT-2 Small. Extension to claims about larger models may reveal that some criteria are harder to satisfy.

\paragraph{Inter-rater reliability.} While the 36 criteria are operationally defined, the aggregation from criterion-level results to verdict tiers involves judgment. Two evaluators could assign different tiers to the same circuit if they weight partial evidence differently. The framework mitigates this by design: the structured validity profile---not the tier---is the primary output. Disagreement is isolated to individual criterion judgments, each of which is independently falsifiable, rather than to a holistic gestalt. We have not yet conducted a formal inter-rater reliability study for the sixteen case-study verdicts presented here. A companion paper addresses this gap through both human inter-annotator agreement (planned: three independent raters on a subset of cases) and automated verdict assignment via structured LLM extraction, where the automation serves as a reproducibility check rather than a replacement for expert judgment.

\paragraph{Case studies are illustrative, not pre-registered.} The sixteen verdicts were assigned after the framework was designed, creating a risk of retrofitting---the same failure mode Mechanistic Validity is designed to detect. We acknowledge this explicitly: the case studies demonstrate how the framework operates and show that it produces discriminating verdicts, but they are not independent tests of the framework's calibration. Pre-registered application to new circuits as they are published---where the verdict is committed before the evidence is fully reviewed---would provide stronger validation and is planned as future work.

\paragraph{Tier thresholds are provisional.} The minimum evidence requirements for tier promotion (\Cref{tab:tier-promotion})---such as ``$\geq 3$ evidence families'' for Triangulated or ``$\geq 2$ ablation methods'' for Mechanistically Supported---are not derived from first principles. They are calibrated against the authors' judgment on the case studies and against established precedents in the source disciplines (clinical trials require multiple endpoints; genetics requires genome-wide replication). The synthesis methods in \Cref{sec:methodology}---Robust Rank Aggregation for cross-method convergence, Dawid-Skene consensus for reliability-weighted voting, and distributional stability via Wasserstein distance---partially alleviate this concern by replacing fixed count thresholds with principled statistical aggregation: a ``$\geq 3$ evidence families'' requirement becomes ``statistically significant convergence across independent methods'' when operationalized through RRA. Nonetheless, the mapping from continuous synthesis scores to discrete tier boundaries involves judgment, and a formal sensitivity analysis showing how verdicts change under threshold perturbation would strengthen the framework.

\paragraph{Calibration of the framework itself.} A framework is an instrument, and instruments are judged by calibration rather than by refutation: GRADE, the USPSTF grading scheme and the multitrait--multimethod matrix issue no falsifiable predictions either, and no one has asked what would refute them. The commitment we make instead is to a ground-truth case. Modular addition supplies two rival algorithms known independently of any audit, and Fourier evidence does not discriminate them---on a Pizza model the Clock functional form still explains 75.4--75.6\% of logit variance \citep[Table~2]{zhong2023clock}. Running the measurement and interpretive suites there should certify the algorithm the model implements and deny that tier to the other, on the strength of the discriminating instruments rather than the shared one. A framework that certifies both, or neither, is miscalibrated. That test can be run today, which is what makes it a commitment rather than a disclaimer.

\paragraph{Scope of ``method-agnostic.''}  The framework is applied almost entirely to attention-head-level circuit claims in GPT-2-scale models. Extending it to SAE features, steering vectors, or crosscoders may require adapting how specific criteria are operationalized---the necessity/sufficiency tests look different for a distributed feature than for a localized head. The criteria themselves (``is specificity tested?'', ``does it generalize?'') are method-agnostic; their concrete instantiation is not. We have gestured at this adaptation in the metrics inventory (\Cref{appendix:metrics}) but have not demonstrated a full end-to-end application to a non-circuit MI result.

\section{Conclusion}
\label{sec:conclusion}

The Mechanistic Validity framework evaluates mechanistic interpretability claims through 36 criteria organized across five validity types in dependency order. Each claim receives a structured validity profile that makes gaps visible and the path to strengthening them concrete.

The framework is open-source and method-agnostic. It applies the same criteria to circuits, SAE features, probing classifiers, steering vectors, transcoders, and crosscoders---because the questions ``Is the construct well-defined?'' and ``Is the evidence causal?'' do not depend on how the evidence was produced.

% ============================================================
% \section*{Acknowledgments}
%
% The theoretical framework draws on insights from the mechanistic interpretability community, particularly the open problems identified by \citet{sharkey2025open} and the actionability framework proposed by \citet{orgad2026actionable}. The cross-field methodology benefits from established validation frameworks in philosophy of science \citep{mayo2018statistical}, neuroscience, psychometrics, pharmacology, and causal inference \citep{pearl2009causality, spirtes2000causation}.
%
% ============================================================

\bibliographystyle{tmlr}

\nocite{campbell1963experimental,loevinger1957objective}
\bibliography{references}


\clearpage
\appendix
\section{Appendix}

This appendix provides the full metric inventory grouped by evidence family (\ref{appendix:metrics}).

\subsection{Metrics inventory}
\label{appendix:metrics}
The full inventory of metrics is listed below, grouped by what each metric measures: causal, structural, information-theoretic, behavioral, representational, and measurement. These groupings organize the catalog. The evidence families of \Cref{sec:evidence-families} classify by source of signal instead, so a metric group may draw on more than one family. \\
Groups A--E are candidate metrics drawn from adjacent literatures: they are available to the framework and most are not exercised by the audits in this paper.  Group F is different in kind.  Those fifteen are the calibrations the audits actually run, they are the only group that maps onto the criteria (\Cref{tab:f-to-m}), and a claim's measurement verdict is decided by them.
\\
\begingroup
\small
\setlength{\tabcolsep}{4pt}
\renewcommand{\arraystretch}{1.15}

% --- A. Causal ---

\begingroup
\small
\setlength{\tabcolsep}{4pt}
\renewcommand{\arraystretch}{1.15}
\noindent\textbf{A. Causal}\par
\vspace{4pt}\noindent\captionof{table}{Causal metrics: the effect of interventions on the model's computation.}
\nopagebreak

\centering
\begin{tabular}{@{}lp{10cm}@{}}
\toprule
\textbf{Metric} & \textbf{Description} \\
\midrule
A01 Counterfactual DAS & Distributed Alignment Search and IIA for testing causal abstraction hypotheses \\
\midrule
A02 Rubin CATE & Conditional average treatment effects across input subpopulations \\
\midrule
A03 Mediation & Natural direct and indirect effects decomposition for circuit paths \\
\midrule
A04 MDL / SLT & Minimum description length and singular learning theory for circuit complexity \\
\midrule
A05 Actual Cause & Halpern-Pearl actual causation on specific inputs, not just potential causes \\
\midrule
A06 Transportability & Formal conditions for circuit generalization across models and distributions \\
\midrule
A07 Causal Discovery & NOTEARS / PC algorithms for learning DAGs over circuit components \\
\midrule
A08 Onset--Offset Coupling & Whether a mechanism's measured strength tracks the capability in both directions: rising as the capability is acquired during training, falling as it is removed by unlearning or fine-tuning, and returning if the capability returns. Adapted from guideline 3 of the molecular Koch's postulates \citep{fredricks1996sequence} \\
\bottomrule
\end{tabular}
\medskip



\begin{minipage}{\linewidth}
\noindent{\textbf{B. Structural}}\par

\begin{center}
\captionof{table}{Structural metrics: weight-space properties requiring no forward pass.}

\begin{tabular}{@{}lp{10cm}@{}}
\toprule
\textbf{Metric} & \textbf{Description} \\
\midrule
B01 SVD / Spectral & Singular value decomposition to identify dominant computational directions \\
\midrule
B02 Effective Rank & Entropy-based dimensionality as a scalar summary of spectral concentration \\
\midrule
B03 OV / QK Decomp. & Decomposing attention heads into OV (what to write) and QK (where to attend) \\
\midrule
B04 Weight Alignment & Cosine similarity between principal weight directions across heads \\
\midrule
B05 Norm Trajectory & Spectral norm ratios tracking signal amplification through components \\
\midrule
B06 Template Distance & Graph-edit and metric distances between circuits discovered for different tasks \\
\midrule
B07 Polysemanticity & Measuring whether components encode multiple unrelated features in superposition \\
\midrule
B08 ICA / NMF & Independent component analysis for decomposing weights into interpretable parts \\
\midrule
B09 Weight Classifier & Training classifiers on weight matrices to predict circuit membership \\
\bottomrule
\end{tabular}
\end{center}
\end{minipage}

\smallskip

% --- information flow and dependencies.
\begin{minipage}{\linewidth}
\noindent{\textbf{C. Information-theoretic}}\par
\begin{center}
\captionof{table}{Information-theoretic metrics: information flow and dependencies.}

\begin{tabular}{@{}lp{10cm}@{}}
\toprule
\textbf{Metric} & \textbf{Description} \\
\midrule
C01 Mutual Info. & Total shared information between circuit components and task performance \\
\midrule
C02 Conditional MI & Information shared after conditioning on other parts of the circuit \\
\midrule
C03 Transfer Entropy & Directed information flow between components across layers \\
\midrule
C04 PID & Decomposing shared information into unique, redundant, and synergistic atoms \\
\midrule
C05 Info. Bottleneck & How efficiently circuits compress input while preserving task-relevant signal \\
\midrule
C06 O-Information & Whether a group of components interacts redundantly or synergistically \\
\midrule
C07 Granger Causality & Whether one component's past activations improve prediction of another's future \\
\midrule
C08 OCSE & Estimating causal influence between components using only observational data \\
\midrule
\bottomrule
\end{tabular}
\end{center}
\end{minipage}

% --- input-output behavior under manipulation.
\begin{minipage}{\linewidth}
\noindent{\textbf{D. Behavioral}}\par
\begin{center}
\captionof{table}{Behavioral metrics: input-output behavior under manipulation.}
\begin{tabular}{@{}lp{10cm}@{}}
\toprule
\textbf{Metric} & \textbf{Description} \\
\midrule
D01 Faithfulness & Whether an identified circuit faithfully reproduces the full model's behavior \\
\midrule
D02 Logit Diff & How much of the model's logit difference the circuit recovers \\
\midrule
D03 KL Divergence & Information-theoretic distance between circuit and full model distributions \\
\midrule
D04 CE Delta & Change in cross-entropy loss when the circuit is ablated \\
\midrule
D05 Top-K Accuracy & Whether the circuit preserves the model's top-K predicted tokens \\
\midrule
D06 Cross-Task & Whether a circuit discovered on one task transfers to a different task \\
\midrule
D07 Cross-Scale & Whether circuit structure replicates in larger or smaller models \\
\midrule
D08 Prompt Paraphrase & Circuit consistency across semantically equivalent prompt templates \\
\midrule
D09 Generalization Gap & Sensitivity of circuit discovery to hyperparameters and methodological choices \\
\bottomrule
\end{tabular}
\end{center}
\end{minipage}

\smallskip

% --- geometry and content of internal representations
\begin{minipage}{\linewidth}
\noindent{\textbf{E. Representational}}\par
\begin{center}
\captionof{table}{Representational metrics: geometry and content of internal representations.}
\begin{tabular}{@{}lp{10cm}@{}}
\toprule
\textbf{Metric} & \textbf{Description} \\
\midrule
E01 DAS-IIA & Whether a learned linear subspace causally encodes a target variable \\
\midrule
E02 Linear Probe & Whether a target variable is linearly decodable from intermediate representations \\
\midrule
E03 RSA & Comparing representation geometry by correlating pairwise distance matrices \\
\midrule
E04 CKA & Kernel alignment for cross-layer and cross-model representation comparison \\
\midrule
E05 Subspace Align. & Cosine alignment between SVD-derived principal directions of weight matrices \\
\midrule
E06 PCA Dim. & Effective dimensionality of circuit subspaces via activation covariance spectrum \\
\midrule
E07 Intrinsic Dim. & True manifold dimensionality, connecting to geometric complexity \\
\midrule
E08 Participation Ratio & How many dimensions are effectively active in a representation \\
\midrule
E09 Persistent Homology & Topological data analysis detecting loops and voids in activation manifolds \\
\midrule
E10 Cross-Task Overlap & Representational structure shared between tasks via IIA transfer \\
\bottomrule
\end{tabular}
\end{center}
\end{minipage}

% --- calibration checks for the trustworthiness of other metrics.

\begin{minipage}{\linewidth}
\noindent{\textbf{F. Measurement}}\par
\begin{center}
\captionof{table}{Measurement metrics: calibration checks for the trustworthiness of other metrics.}
\begin{tabular}{@{}lp{10cm}@{}}
\toprule
\textbf{Metric} & \textbf{Description} \\
\midrule
F01 Test--Retest & Whether rerunning the same metric on the same model, task, and intervention gives the same answer \\
\midrule
F02 Bootstrap Stability & Whether the metric remains stable under bootstrap resampling of prompts, examples, or activation samples \\
\midrule
F03 Seed Variance & How much the metric changes across random seeds for probes, SAEs, optimization, or sampling-based estimators \\
\midrule
F04 Checkpoint Variance & Whether the metric is stable across nearby checkpoints or depends on one training snapshot \\
\midrule
F05 Prompt Variance & How much the metric changes across prompt templates, paraphrases, lexical choices, or dataset slices \\
\midrule
F06 Baseline Separation & Whether the score is distinguishable from random, untrained, shuffled-label, or permuted-circuit baselines \\
\midrule
F07 Negative Controls & Whether the metric stays low where the claimed effect should be absent \\
\midrule
F08 Positive Controls & Whether the metric detects known-true or synthetic effects of the relevant size \\
\midrule
F09 Intervention Robustness & Whether the conclusion survives reasonable intervention variants such as mean ablation, zero ablation, resample ablation, patching, or causal scrubbing \\
\midrule
F10 Hyperparam. Sensitivity & Whether the metric changes under reasonable choices of sparsity, probe regularization, thresholds, localization cutoffs, or discovery hyperparameters \\
\midrule
F11 Estimator Uncertainty & Confidence intervals, standard errors, permutation tests, or posterior intervals for the reported metric value \\
\midrule
F12 Calibration Curve & Whether larger reported scores correspond to larger empirical effects or higher probability of success \\
\midrule
F13 Cross-Metric Convergence & Whether independent metrics intended to test the same criterion agree despite different failure modes \\
\midrule
F14 Measurement Invariance & Whether the metric behaves consistently across model sizes, tasks, prompt distributions, and intervention contexts \\
\midrule
F15 Multiple Comparisons & Whether scores are corrected for the number of components tested, methods tried, or thresholds swept. Reporting the top-$k$ heads out of $N$ without FDR or permutation correction inflates apparent circuit specificity \\
\bottomrule
\end{tabular}
\end{center}
\end{minipage}

\endgroup

\end{document}
